%%
% Copyright (c) 2017 - 2026, Pascal Wagler;
% Copyright (c) 2014 - 2026, John MacFarlane
%
% All rights reserved.
%
% Redistribution and use in source and binary forms, with or without
% modification, are permitted provided that the following conditions
% are met:
%
% - Redistributions of source code must retain the above copyright
% notice, this list of conditions and the following disclaimer.
%
% - Redistributions in binary form must reproduce the above copyright
% notice, this list of conditions and the following disclaimer in the
% documentation and/or other materials provided with the distribution.
%
% - Neither the name of John MacFarlane nor the names of other
% contributors may be used to endorse or promote products derived
% from this software without specific prior written permission.
%
% THIS SOFTWARE IS PROVIDED BY THE COPYRIGHT HOLDERS AND CONTRIBUTORS
% "AS IS" AND ANY EXPRESS OR IMPLIED WARRANTIES, INCLUDING, BUT NOT
% LIMITED TO, THE IMPLIED WARRANTIES OF MERCHANTABILITY AND FITNESS
% FOR A PARTICULAR PURPOSE ARE DISCLAIMED. IN NO EVENT SHALL THE
% COPYRIGHT OWNER OR CONTRIBUTORS BE LIABLE FOR ANY DIRECT, INDIRECT,
% INCIDENTAL, SPECIAL, EXEMPLARY, OR CONSEQUENTIAL DAMAGES (INCLUDING,
% BUT NOT LIMITED TO, PROCUREMENT OF SUBSTITUTE GOODS OR SERVICES;
% LOSS OF USE, DATA, OR PROFITS; OR BUSINESS INTERRUPTION) HOWEVER
% CAUSED AND ON ANY THEORY OF LIABILITY, WHETHER IN CONTRACT, STRICT
% LIABILITY, OR TORT (INCLUDING NEGLIGENCE OR OTHERWISE) ARISING IN
% ANY WAY OUT OF THE USE OF THIS SOFTWARE, EVEN IF ADVISED OF THE
% POSSIBILITY OF SUCH DAMAGE.
%%

%%
% This is the Eisvogel pandoc LaTeX template.
%
% For usage information and examples visit the official GitHub page:
% https://github.com/Wandmalfarbe/pandoc-latex-template
%%
% Options for packages loaded elsewhere
\PassOptionsToPackage{unicode}{hyperref}
\PassOptionsToPackage{hyphens}{url}
\PassOptionsToPackage{dvipsnames,svgnames,x11names,table}{xcolor}
\documentclass[
  spanish,
  11pt,
  paper=a4,
  ,captions=tableheading
]{scrartcl}
\usepackage{xcolor}
\usepackage[margin=2.5cm]{geometry}
\usepackage{amsmath,amssymb}
\usepackage[export]{adjustbox}
\usepackage{graphicx}

% add backlinks to footnote references, cf. https://tex.stackexchange.com/questions/302266/make-footnote-clickable-both-ways
\usepackage{footnotebackref}
\setcounter{secnumdepth}{-\maxdimen} % remove section numbering
\usepackage{iftex}
\ifPDFTeX
  \usepackage[T1]{fontenc}
  \usepackage[utf8]{inputenc}
  \usepackage{textcomp} % provide euro and other symbols
\else % if luatex or xetex
  \usepackage{unicode-math} % this also loads fontspec
  \defaultfontfeatures{Scale=MatchLowercase}
  \defaultfontfeatures[\rmfamily]{Ligatures=TeX,Scale=1}
\fi
\usepackage{lmodern}
\ifPDFTeX\else
  % xetex/luatex font selection
\fi
% Use upquote if available, for straight quotes in verbatim environments
\IfFileExists{upquote.sty}{\usepackage{upquote}}{}
\IfFileExists{microtype.sty}{% use microtype if available
  \usepackage[]{microtype}
  \UseMicrotypeSet[protrusion]{basicmath} % disable protrusion for tt fonts
}{}

% Use setspace anyway because we change the default line spacing.
% The spacing is changed early to affect the titlepage and the TOC.
\usepackage{setspace}
\setstretch{1.2}
\makeatletter
\@ifundefined{KOMAClassName}{% if non-KOMA class
  \IfFileExists{parskip.sty}{%
    \usepackage{parskip}
  }{% else
    \setlength{\parindent}{0pt}
    \setlength{\parskip}{6pt plus 2pt minus 1pt}}
}{% if KOMA class
  \KOMAoptions{parskip=half}}
\makeatother
\usepackage{color}
\usepackage{fancyvrb}
\newcommand{\VerbBar}{|}
\newcommand{\VERB}{\Verb[commandchars=\\\{\}]}
\DefineVerbatimEnvironment{Highlighting}{Verbatim}{commandchars=\\\{\}}
% Add ',fontsize=\small' for more characters per line
\newenvironment{Shaded}{}{}
\newcommand{\AlertTok}[1]{\textcolor[rgb]{1.00,0.00,0.00}{\textbf{#1}}}
\newcommand{\AnnotationTok}[1]{\textcolor[rgb]{0.38,0.63,0.69}{\textbf{\textit{#1}}}}
\newcommand{\AttributeTok}[1]{\textcolor[rgb]{0.49,0.56,0.16}{#1}}
\newcommand{\BaseNTok}[1]{\textcolor[rgb]{0.25,0.63,0.44}{#1}}
\newcommand{\BuiltInTok}[1]{\textcolor[rgb]{0.00,0.50,0.00}{#1}}
\newcommand{\CharTok}[1]{\textcolor[rgb]{0.25,0.44,0.63}{#1}}
\newcommand{\CommentTok}[1]{\textcolor[rgb]{0.38,0.63,0.69}{\textit{#1}}}
\newcommand{\CommentVarTok}[1]{\textcolor[rgb]{0.38,0.63,0.69}{\textbf{\textit{#1}}}}
\newcommand{\ConstantTok}[1]{\textcolor[rgb]{0.53,0.00,0.00}{#1}}
\newcommand{\ControlFlowTok}[1]{\textcolor[rgb]{0.00,0.44,0.13}{\textbf{#1}}}
\newcommand{\DataTypeTok}[1]{\textcolor[rgb]{0.56,0.13,0.00}{#1}}
\newcommand{\DecValTok}[1]{\textcolor[rgb]{0.25,0.63,0.44}{#1}}
\newcommand{\DocumentationTok}[1]{\textcolor[rgb]{0.73,0.13,0.13}{\textit{#1}}}
\newcommand{\ErrorTok}[1]{\textcolor[rgb]{1.00,0.00,0.00}{\textbf{#1}}}
\newcommand{\ExtensionTok}[1]{#1}
\newcommand{\FloatTok}[1]{\textcolor[rgb]{0.25,0.63,0.44}{#1}}
\newcommand{\FunctionTok}[1]{\textcolor[rgb]{0.02,0.16,0.49}{#1}}
\newcommand{\ImportTok}[1]{\textcolor[rgb]{0.00,0.50,0.00}{\textbf{#1}}}
\newcommand{\InformationTok}[1]{\textcolor[rgb]{0.38,0.63,0.69}{\textbf{\textit{#1}}}}
\newcommand{\KeywordTok}[1]{\textcolor[rgb]{0.00,0.44,0.13}{\textbf{#1}}}
\newcommand{\NormalTok}[1]{#1}
\newcommand{\OperatorTok}[1]{\textcolor[rgb]{0.40,0.40,0.40}{#1}}
\newcommand{\OtherTok}[1]{\textcolor[rgb]{0.00,0.44,0.13}{#1}}
\newcommand{\PreprocessorTok}[1]{\textcolor[rgb]{0.74,0.48,0.00}{#1}}
\newcommand{\RegionMarkerTok}[1]{#1}
\newcommand{\SpecialCharTok}[1]{\textcolor[rgb]{0.25,0.44,0.63}{#1}}
\newcommand{\SpecialStringTok}[1]{\textcolor[rgb]{0.73,0.40,0.53}{#1}}
\newcommand{\StringTok}[1]{\textcolor[rgb]{0.25,0.44,0.63}{#1}}
\newcommand{\VariableTok}[1]{\textcolor[rgb]{0.10,0.09,0.49}{#1}}
\newcommand{\VerbatimStringTok}[1]{\textcolor[rgb]{0.25,0.44,0.63}{#1}}
\newcommand{\WarningTok}[1]{\textcolor[rgb]{0.38,0.63,0.69}{\textbf{\textit{#1}}}}

% Workaround/bugfix from jannick0.
% See https://github.com/jgm/pandoc/issues/4302#issuecomment-360669013)
% or https://github.com/Wandmalfarbe/pandoc-latex-template/issues/2
%
% Redefine the verbatim environment 'Highlighting' to break long lines (with
% the help of fvextra). Redefinition is necessary because it is unlikely that
% pandoc includes fvextra in the default template.
\usepackage{fvextra}
\DefineVerbatimEnvironment{Highlighting}{Verbatim}{breaklines,fontsize=\small,commandchars=\\\{\}}

\usepackage{longtable,booktabs,array}
\newcounter{none} % for unnumbered tables
\usepackage{calc} % for calculating minipage widths
% Correct order of tables after \paragraph or \subparagraph
\usepackage{etoolbox}
\makeatletter
\patchcmd\longtable{\par}{\if@noskipsec\mbox{}\fi\par}{}{}
\makeatother
% Allow footnotes in longtable head/foot
\IfFileExists{footnotehyper.sty}{\usepackage{footnotehyper}}{\usepackage{footnote}}
\makesavenoteenv{longtable}
\usepackage{graphicx}
\makeatletter
\newsavebox\pandoc@box
\newcommand*\pandocbounded[1]{% scales image to fit in text height/width
  \sbox\pandoc@box{#1}%
  \Gscale@div\@tempa{\textheight}{\dimexpr\ht\pandoc@box+\dp\pandoc@box\relax}%
  \Gscale@div\@tempb{\linewidth}{\wd\pandoc@box}%
  \ifdim\@tempb\p@<\@tempa\p@\let\@tempa\@tempb\fi% select the smaller of both
  \ifdim\@tempa\p@<\p@\scalebox{\@tempa}{\usebox\pandoc@box}%
  \else\usebox{\pandoc@box}%
  \fi%
}
% Set default figure placement to htbp
% Make use of float-package and set default placement for figures to H.
% The option H means 'PUT IT HERE' (as  opposed to the standard h option which means 'You may put it here if you like').
\usepackage{float}
\floatplacement{figure}{H}
\makeatother
\ifLuaTeX
\usepackage[bidi=basic,shorthands=off]{babel}
\else
\usepackage[bidi=default,shorthands=off]{babel}
\fi
\ifLuaTeX
  \usepackage{selnolig} % disable illegal ligatures
\fi
\setlength{\emergencystretch}{3em} % prevent overfull lines
\providecommand{\tightlist}{%
  \setlength{\itemsep}{0pt}\setlength{\parskip}{0pt}}
\usepackage{xurl}
\usepackage{bookmark}
\IfFileExists{xurl.sty}{\usepackage{xurl}}{} % add URL line breaks if available
\urlstyle{same}
\definecolor{default-linkcolor}{HTML}{A50000}
\definecolor{default-filecolor}{HTML}{A50000}
\definecolor{default-citecolor}{HTML}{4077C0}
\definecolor{default-urlcolor}{HTML}{4077C0}

\hypersetup{
  pdftitle={Informe del Proyecto: Agente Gamer},
  pdfauthor={Javier Lesaca Medina},
  pdflang={es},
  hidelinks,
  breaklinks=true,
  pdfcreator={LaTeX via pandoc with the Eisvogel template}}

\title{Informe del Proyecto: Agente Gamer}
\author{Javier Lesaca Medina}
\date{Curso 2025/26}


%
% for the background color of the title page
%
\usepackage{pagecolor}
\usepackage{afterpage}

%
% break urls
%
\PassOptionsToPackage{hyphens}{url}

%
% When using babel or polyglossia with biblatex, loading csquotes is recommended
% to ensure that quoted texts are typeset according to the rules of your main language.
%
\usepackage{csquotes}

%
% captions
%
\definecolor{caption-color}{HTML}{777777}
\usepackage[font={stretch=1.2}, textfont={color=caption-color}, position=top, skip=4mm, labelfont=bf, singlelinecheck=false, justification=raggedright]{caption}
\setcapindent{0em}

%
% blockquote
%
\definecolor{blockquote-border}{RGB}{221,221,221}
\definecolor{blockquote-text}{RGB}{119,119,119}
\usepackage{mdframed}
\newmdenv[rightline=false,bottomline=false,topline=false,linewidth=3pt,linecolor=blockquote-border,skipabove=\parskip]{customblockquote}
\renewenvironment{quote}{\begin{customblockquote}\list{}{\rightmargin=0em\leftmargin=0em}%
\item\relax\color{blockquote-text}\ignorespaces}{\unskip\unskip\endlist\end{customblockquote}}

%
% Source Sans Pro as the default font family
% Source Code Pro for monospace text
%
% 'default' option sets the default
% font family to Source Sans Pro, not \sfdefault.
%
% Note that the font has been officially renamed to `Source Sans 3`, and
% the version provided by the `sourcesanspro` package is slightly outdated.
% You can install the newer version locally and use it, for example, with
% `mainfont: "Source Sans 3"` in the YAML metadata (requires XeTeX or LuaTeX).
%
\ifnum 0\ifxetex 1\fi\ifluatex 1\fi=0 % if pdftex
    \usepackage[default]{sourcesanspro}
  \usepackage{sourcecodepro}
  \else % if not pdftex
    \usepackage[default]{sourcesanspro}
  \usepackage{sourcecodepro}

  % XeLaTeX specific adjustments for straight quotes: https://tex.stackexchange.com/a/354887
  % This issue is already fixed (see https://github.com/silkeh/latex-sourcecodepro/pull/5) but the
  % fix is still unreleased.
  % TODO: Remove this workaround when the new version of sourcecodepro is released on CTAN.
  \ifxetex
    \makeatletter
    \defaultfontfeatures[\ttfamily]
      { Numbers   = \sourcecodepro@figurestyle,
        Scale     = \SourceCodePro@scale,
        Extension = .otf }
    \setmonofont
      [ UprightFont    = *-\sourcecodepro@regstyle,
        ItalicFont     = *-\sourcecodepro@regstyle It,
        BoldFont       = *-\sourcecodepro@boldstyle,
        BoldItalicFont = *-\sourcecodepro@boldstyle It ]
      {SourceCodePro}
    \makeatother
  \fi
  \fi

%
% heading color
%
\definecolor{heading-color}{RGB}{40,40,40}
% By default, the KOMA-Script classes will typeset sectioning headings in
% sans-serif. Use the normal body font for headings.
\addtokomafont{disposition}{\normalfont\color{heading-color}\bfseries}

%
% variables for title, author and date
%
\usepackage{titling}
\title{Informe del Proyecto: Agente Gamer}
\author{Javier Lesaca Medina}
\date{Curso 2025/26}

%
% tables
%

\definecolor{table-row-color}{HTML}{F5F5F5}
\definecolor{table-rule-color}{HTML}{999999}

%\arrayrulecolor{black!40}
\arrayrulecolor{table-rule-color}     % color of \toprule, \midrule, \bottomrule
\setlength\heavyrulewidth{0.3ex}      % thickness of \toprule, \bottomrule
\renewcommand{\arraystretch}{1.3}     % spacing (padding)


%
% remove paragraph indentation
%
\setlength{\parindent}{0pt}
\setlength{\parskip}{6pt plus 2pt minus 1pt}
\setlength{\emergencystretch}{3em}  % prevent overfull lines

%
%
% Listings
%
%


%
% header and footer
%
\usepackage[headsepline,footsepline]{scrlayer-scrpage}

\newpairofpagestyles{eisvogel-header-footer}{
  \clearpairofpagestyles
  \ihead*{Informe del Proyecto: Agente Gamer}
  \chead*{}
  \ohead*{Curso 2025/26}
  \ifoot*{Javier Lesaca Medina}
  \cfoot*{}
  \ofoot*{\thepage}
  \addtokomafont{pageheadfoot}{\upshape}
}
\pagestyle{eisvogel-header-footer}



%
% Define watermark
%

\begin{document}

\begin{titlepage}
\newgeometry{left=6cm}
\definecolor{titlepage-color}{HTML}{1D689D}
\newpagecolor{titlepage-color}\afterpage{\restorepagecolor}
\newcommand{\colorRule}[3][black]{\textcolor[HTML]{#1}{\rule{#2}{#3}}}
\begin{flushleft}
\noindent
\\[-1em]
\color[HTML]{FFFFFF}
\makebox[0pt][l]{\colorRule[FFFFFF]{1.3\textwidth}{2pt}}
\par
\noindent

{
  \setstretch{1.4}
  \vfill
  \noindent {\huge \textbf{\textsf{Informe del Proyecto: Agente Gamer}}}
    \vskip 2em
  \noindent {\Large \textsf{Javier Lesaca Medina}}
  \vfill
}

\noindent
\includegraphics[width=250pt, left]{/home/javier/proyectos/AgenteGamer/AgenteGamer\_APP/docs/img/logo\_safa\_blanco.png}

\textsf{Curso 2025/26}
\end{flushleft}
\end{titlepage}
\restoregeometry
\pagenumbering{arabic}

% don't generate the default title
% \maketitle


{
\setcounter{tocdepth}{3}
\tableofcontents
}
\begin{figure}
\centering

\includegraphics[width=0.7\textwidth,height=\textheight]{docs/img/splash.png}
\caption{Splash‑view de la app}
\end{figure}

\textbf{AGENTE\,GAMER} \emph{``Juega más, gasta mejor''}

\begin{center}\rule{0.5\linewidth}{0.5pt}\end{center}

\hypertarget{agradecimientos}{%
\subsection{Agradecimientos}\label{agradecimientos}}

Quiero expresar mi más sincero agradecimiento a todas las personas que
me acompañaron en este camino:

\begin{itemize}
\tightlist
\item
  \textbf{Mentores y docentes} que compartieron su conocimiento y
  paciencia.
\item
  \textbf{Compañeros y amigos} que probaron la aplicación, detectaron
  bugs y ofrecieron ideas.
\item
  \textbf{Familia} por su apoyo incondicional y por creer en mis
  proyectos.
\end{itemize}

Este trabajo es el resultado de su guía, su tiempo y su entusiasmo.
¡Gracias por hacer posible este proyecto!

\begin{center}\rule{0.5\linewidth}{0.5pt}\end{center}

\hypertarget{cita-motivadora}{%
\subsection{Cita motivadora}\label{cita-motivadora}}

\begin{quote}
\emph{``La mejor manera de predecir el futuro es implementarlo.''}
---\,\textbf{David Heinemeier Hansson}, creador de Ruby on Rails
\end{quote}

\usepackage{xcolor}
  \usepackage{sectsty}
  \definecolor{headingblue}{HTML}{1D689D}
  \allsectionsfont{\color{headingblue}}
  \usepackage{hyperref}
  \usepackage{pdflscape}
  \usepackage{graphicx}
  \usepackage{float}
  \usepackage{placeins}
  \let\origfigure\figure
  \let\endorigfigure\endfigure
  \renewenvironment{figure}[1][2]{\origfigure[H]}{\endorigfigure}
  \hypersetup{
    colorlinks=true,
    linkcolor=headingblue,
    urlcolor=headingblue,
    citecolor=headingblue
  }

\begin{center}\rule{0.5\linewidth}{0.5pt}\end{center}

\newpage

\hypertarget{descripciuxf3n}{%
\section{1.- Descripción}\label{descripciuxf3n}}

En la actualidad, los videojuegos se han consolidado como una de las
formas de entretenimiento más populares a nivel mundial. La industria
del gaming genera miles de millones de euros anuales, con un crecimiento
constante año tras año. Sin embargo, uno de los problemas que enfrentan
muchos jugadores, especialmente los más activos, es la gestión adecuada
de su presupuesto destinado a la adquisición de juegos, suscripciones,
contenidos descargables y otros elementos relacionados con el mundo
gamer.

\textbf{AgenteGamer} es una aplicación Android diseñada específicamente
para ayudar a los jugadores a gestionar de manera inteligente su
presupuesto mensual dedicado a los videojuegos. El nombre de la
aplicación surge del concepto de un «agente» financiero personal que
asesora al usuario sobre cómo optimizar sus gastos en gaming.

El problema que resuelve AgenteGamer es doble: por un lado, muchos
jugadores no tienen visibilidad clara de cuánto dinero gastan
mensualmente en videojuegos; por otro, carecen de herramientas que les
permitan tomar decisiones informadas sobre qué juegos comprar
considerando su situación financiera personal.

La aplicación está dirigida principalmente a:

\begin{itemize}
\tightlist
\item
  \textbf{Jugadores habituales} de entre 16 y 45 años que gastan
  regularmente en videojuegos y buscan controlar su presupuesto.
\item
  \textbf{Estudiantes y jóvenes adultos} que desean administrar de forma
  responsable su presupuesto limitado asignado al gaming.
\item
  \textbf{Coleccionistas y entusiastas} que compran numerosos juegos y
  necesitan herramientas para priorizar compras.
\item
  \textbf{Familias} que establecen presupuestos de entretenimiento para
  sus hijos y desean supervisar el gasto en videojuegos.
\end{itemize}

El \textbf{tagline} de la aplicación es «Juega más, gasta mejor»,
reflejando la filosofía de maximizar el entretenimiento gamer sin
comprometer la estabilidad financiera personal.

La propuesta de valor principal de AgenteGamer radica en combinar tres
funcionalidades que no se encuentran de forma integrada en otras
aplicaciones del mercado: gestión de gastos en videojuegos, catálogo de
juegos con información actualizada, y un sistema de recomendaciones
personalizado que evalúa la conveniencia de compras basándose en el
presupuesto y preferencias del usuario.

\newpage

\hypertarget{anuxe1lisis-del-mercado-actual}{%
\section{2.- Análisis del Mercado
Actual}\label{anuxe1lisis-del-mercado-actual}}

\hypertarget{ynab-you-need-a-budget}{%
\subsection{YNAB (You Need A Budget)}\label{ynab-you-need-a-budget}}

\begin{center}

\includegraphics[width=0.5\textwidth]{/home/javier/proyectos/AgenteGamer/AgenteGamer_APP/docs/img/YNAB-logo.png}
\end{center}

YNAB es una de las aplicaciones de presupuesto personal más populares
del mercado. Su metodología se basa en cuatro reglas fundamentales para
gestionar el dinero. Entre sus puntos fuertes destaca su enfoque
educativo y su comunidad activa. Sin embargo, su modelo de suscripción
(84€/año) resulta elevado para usuarios jóvenes, y no ofrece ninguna
funcionalidad específica para gamers. Su interfaz, aunque funcional, es
genérica y no conecta emocionalmente con el público gaming.

\hypertarget{spendee}{%
\subsection{Spendee}\label{spendee}}

\begin{center}

\includegraphics[width=0.5\textwidth]{/home/javier/proyectos/AgenteGamer/AgenteGamer_APP/docs/img/Spendee_logo.png}
\end{center}

Spendee ofrece una interfaz visual atractiva con gráficos claros y la
posibilidad de compartir presupuestos familiares. Su punto fuerte es la
visualización de datos y la categorización automática de gastos. Por
contra, la versión gratuita es muy limitada, no permite exportar datos
y, al igual que YNAB, no tiene ninguna funcionalidad orientada al mundo
gaming.

\hypertarget{wallet-by-budgetbakers}{%
\subsection{Wallet by BudgetBakers}\label{wallet-by-budgetbakers}}

\begin{center}

\includegraphics[width=0.5\textwidth]{/home/javier/proyectos/AgenteGamer/AgenteGamer_APP/docs/img/wallet_logo.png}
\end{center}

Wallet es una aplicación completa que permite conectar cuentas bancarias
reales para un seguimiento automático de gastos. Ofrece informes
detallados y planificación a largo plazo. Su debilidad principal es la
complejidad de configuración inicial y la dependencia de la conexión
bancaria (no disponible en todos los países). Tampoco tiene enfoque
gaming.

\hypertarget{splitwise}{%
\subsection{Splitwise}\label{splitwise}}

\begin{center}

\includegraphics[width=0.5\textwidth]{/home/javier/proyectos/AgenteGamer/AgenteGamer_APP/docs/img/splitwise_logo.png}
\end{center}

Aunque no es una app de presupuesto personal, Splitwise es muy popular
entre gamers que comparten gastos (suscripciones familiares, juegos
compartidos). Su fortaleza es la gestión de gastos compartidos, pero no
ofrece control de presupuesto individual ni recomendaciones financieras.

\hypertarget{plataformas-web-especializadas}{%
\subsection{Plataformas web
especializadas}\label{plataformas-web-especializadas}}

Además de las aplicaciones de gestión financiera, existen plataformas
web especializadas en el seguimiento de precios de videojuegos que los
jugadores utilizan como referencia para sus compras:

\textbf{SteamDB}: Base de datos no oficial de Steam con precios
históricos y alertas de ofertas. Su fortaleza es la precisión de datos y
el seguimiento técnico de la plataforma. Sin embargo, carece de gestión
de presupuesto, no permite registrar gastos y está limitada
exclusivamente a Steam.

\textbf{Deku Deals}: Agregador de precios multiplataforma (Steam,
PlayStation, Nintendo, Xbox, GOG) con wishlist centralizada y alertas
personalizadas. Su ventaja es la cobertura amplia de tiendas digitales.
No ofrece evaluación financiera ni herramientas de presupuesto,
enfocándose solo en precios y descuentos.

\textbf{PriceCharting}: Rastreador de precios orientado al coleccionismo
de juegos físicos y retro. Proporciona gráficos de evolución de precios
y alertas configurables. Su debilidad principal es el enfoque en
coleccionismo más que en gestión de gastos, sin integración con
catálogos actuales ni recomendaciones financieras.

Estas herramientas, aunque útiles para encontrar ofertas, no sustituyen
una solución integral de gestión financiera gamer, ya que ninguna
combina el seguimiento de precios con el control de presupuesto y
recomendaciones personalizadas que ofrece AgenteGamer.

\newpage

\hypertarget{propuesta-innovadora}{%
\section{3.- Propuesta Innovadora}\label{propuesta-innovadora}}

\hypertarget{punto-diferenciador}{%
\subsection{3.1.- Punto diferenciador}\label{punto-diferenciador}}

El punto diferenciador de Agente Gamer respecto a las aplicaciones
mencionadas es su \textbf{enfoque específico para gamers}. Mientras que
YNAB, Spendee y Wallet son herramientas genéricas de finanzas
personales, Agente Gamer entiende que:

\begin{enumerate}
\def\labelenumi{\arabic{enumi}.}
\item
  \textbf{Los gamers tienen patrones de gasto específicos}: lanzamientos
  de juegos, DLCs, microtransacciones, suscripciones (Game Pass, PS
  Plus, Nintendo Online), y compras impulsivas en ofertas de Steam.
\item
  \textbf{La wishlist es central en la experiencia gaming}: Los gamers
  mantienen listas de juegos que quieren comprar, con precios estimados
  y fechas de lanzamiento. Agente Gamer integra esta funcionalidad
  nativamente, algo que ninguna app de finanzas ofrece. Además, cada
  juego de la wishlist recibe una evaluación financiera automática
  (recomendada/ajustada/no recomendada) basada en el presupuesto y gasto
  actual del mes.
\item
  \textbf{El catálogo de juegos integrado}: La conexión con RAWG API
  permite buscar juegos, ver precios estimados (calculados según rating)
  y fechas de lanzamiento sin salir de la app de finanzas. La búsqueda
  tiene debounce para optimizar llamadas y paginación infinita para
  mejor experiencia.
\item
  \textbf{Sistema de evaluación financiera inteligente}: El
  SistemaFinanciero evalúa cada compra potencial según umbrales
  porcentuales del presupuesto restante:

  \begin{itemize}
  \item
    \(\leq\) 30\% del restante: ``Compra recomendada''
  \item
    \(\leq\) 60\% del restante: ``Compra ajustada''
  \item
    \(\leq\) 100\% del restante: ``Mejor esperar''
  \item
    \begin{quote}
    100\%: ``No recomendable''
    \end{quote}
  \end{itemize}
\item
  \textbf{Diseño visual gaming}: La interfaz dark theme con acentos
  verdes está diseñada para resonar con la estética gaming, no con la
  estética corporativa de las apps financieras tradicionales.
\item
  \textbf{Notificaciones inteligentes en segundo plano}: WorkManager
  ejecuta workers periódicos que validan la wishlist contra el
  presupuesto y notifican al usuario cuando hay compras no recomendadas
  o lanzamientos próximos (0-7 días).
\item
  \textbf{El ``Agente Financiero''}: Un sistema de recomendaciones
  personalizadas que analiza los patrones de gasto y ofrece consejos
  contextuales, como un asistente virtual que entiende tanto de finanzas
  como de gaming.
\end{enumerate}

Además, la aplicación es \textbf{gratuita}, sin modelo de suscripción,
lo que la hace accesible para el público joven que compone la mayoría de
la comunidad gamer.

\newpage

\hypertarget{funcionalidades-principales}{%
\subsection{3.2.- Funcionalidades
principales}\label{funcionalidades-principales}}

\textbf{Autenticación y Perfil:} - Login y registro con Firebase Auth -
Perfil de usuario sincronizado en Firestore (nombre, email, presupuesto,
moneda, rol, fecha de creación) - Edición de nombre mediante
ActivityResultLauncher - Roles de usuario: USER y ADMIN (los admins
tienen acceso a funciones de prueba)

\textbf{Dashboard Principal (MainActivity):} - Card grande con total
gastado, estado financiero (indicador de color) y recomendación del
agente - Dos cards laterales: presupuesto mensual y restante disponible
- Gráfico de tendencia de gastos mensual (LineChart de MPAndroidChart) -
Indicador de tendencia vs mes anterior (porcentaje con flecha) -
Recomendaciones personalizadas del agente financiero basadas en patrones
de gasto - Lista horizontal de últimos gastos (RecyclerView) - Gráfico
circular de gastos por categoría (PieChart de MPAndroidChart) -
Presupuesto reactivo en tiempo real mediante LiveData

\textbf{Gestión de Gastos (ListaGastosActivity):} - Lista de gastos con
imagen, nombre, precio y fecha - Cálculo y muestra del total de gastos
del período - Empty state cuando no hay gastos - Botón FAB para agregar
gastos de prueba (solo visible para usuarios ADMIN) - Formateo de
totales según la moneda del usuario

\textbf{Catálogo de Juegos (ListaJuegosActivity):} - Grid de 2 columnas
con portadas de juegos obtenidas de la API de RAWG - Barra de búsqueda
con debounce de 500ms para evitar llamadas excesivas - Paginación
infinita al hacer scroll - Precarga de juegos desde SplashActivity para
mejor rendimiento - Al hacer clic en un juego, se agrega a la wishlist
con precio estimado calculado según rating - Moneda reactiva que se
actualiza cuando el usuario cambia en Ajustes

\textbf{Wishlist (ListaWishlistActivity):} - Lista de juegos deseados
con imagen, nombre, precio estimado y evaluación financiera - Evaluación
en tiempo real: ``Compra recomendada'', ``Compra ajustada'', ``Mejor
esperar'', ``No recomendable'' - Cálculo y muestra del costo total
estimado - Diálogo de detalle del juego con edición de precio,
plataforma y confirmación de compra - Al comprar un juego: si es
lanzamiento futuro se guarda en Lanzamientos, si ya está lanzado se crea
un Gasto - Moneda reactiva desde Firestore

\textbf{Lanzamientos (LanzamientosActivity):} - Lista de próximos
lanzamientos (próximos 15 días) - Muestra días restantes hasta el
lanzamiento - Muestra rating del juego y plataformas disponibles -
Precarga de datos desde la API de RAWG al abrir la pantalla

\textbf{Configuración (AjustesActivity):} - Editor de presupuesto
mensual con validación - Selector de moneda (EUR, USD, GBP) con
actualización en tiempo real - Cálculo de impacto en tiempo real basado
en la wishlist actual (recomendados, ajustados, no recomendados) -
Guarda los cambios en Firestore mediante UseCase

\textbf{Perfil (PerfilActivity):} - Muestra email, nombre, presupuesto,
fecha de creación y rol - Presupuesto reactivo en tiempo real - Botón
para editar nombre que lanza EditProfileActivity - Indicador de carga
durante actualización

\newpage

\hypertarget{estudio-de-viabilidad}{%
\section{4.- Estudio de Viabilidad}\label{estudio-de-viabilidad}}

\hypertarget{descripciuxf3n-del-modelo-de-negocio}{%
\subsection{4.1.- Descripción del modelo de
negocio}\label{descripciuxf3n-del-modelo-de-negocio}}

Agente Gamer se plantea inicialmente como un \textbf{proyecto educativo
de código abierto}, sin ánimo de lucro. Sin embargo, si se considerara
su comercialización futura, el modelo de negocio más adecuado sería
\textbf{freemium}:

\begin{itemize}
\tightlist
\item
  \textbf{Plan gratuito}: Gestión básica de gastos, wishlist, catálogo
  de juegos y recomendaciones del agente.
\item
  \textbf{Plan premium} (2,99€/mes): Análisis avanzado de tendencias,
  exportación de datos, múltiples presupuestos, alertas de precios de
  juegos y sincronización multi-dispositivo avanzada.
\end{itemize}

Los costes de mantenimiento incluirían el hosting de Firebase (plan
Spark gratuito para desarrollo, plan Blaze para producción) y el dominio
de la API RAWG.

\hypertarget{embudo-de-marketing}{%
\subsection{4.2.- Embudo de marketing}\label{embudo-de-marketing}}

El embudo de marketing para AgenteGamer se estructura en varias fases
que buscan atraer, convertir y retener usuarios.

\textbf{Fase de atracción (Awareness):} El canal principal será la
optimización para la tienda Google Play con palabras clave relevantes
(«gestión presupuesto gaming», «control gastos videojuegos», «gestor
económico gamer»). Adicionalmente, se utilizarán publicaciones en
comunidades de Reddit relacionadas con gaming (r/gaming,
r/patientgamers, r/Steam), foros especializados y presencia en redes
sociales como Instagram y TikTok con contenido educativo sobre gestión
financiera en gaming.

\textbf{Fase de conversión (Consideration):} Se ofrecerán guías
descargables sobre gestión de presupuesto gaming, calculadora online de
presupuesto gamer en la página web, y prueba gratuita de 7 días de la
versión premium para nuevos usuarios.

\textbf{Fase de retención (Retention):} Notificaciones push
personalizadas con recordatorios mensuales de presupuesto, contenido
exclusivo en redes sociales para usuarios premium, y programa de
referidos con beneficios para ambos usuarios.

\textbf{Fase de monetización:} Llamadas a la acción claras para upgrade
a premium dentro de la aplicación, ofertas especiales en momentos clave
(Black Friday, Navidad), y bundles temáticos (por ejemplo, «Pack Fin de
Año Gamer»).

\newpage

\begin{center}
\includegraphics[width=\textwidth,height=\textheight,keepaspectratio]{/home/javier/proyectos/AgenteGamer/AgenteGamer_APP/docs/img/embudo-marketing.png}
\end{center}

\newpage

\hypertarget{estimaciuxf3n-de-la-poblaciuxf3n-objetivo}{%
\subsection{4.3.- Estimación de la población
objetivo}\label{estimaciuxf3n-de-la-poblaciuxf3n-objetivo}}

El público objetivo son \textbf{gamers de entre 16 y 45 años} que
realizan compras de videojuegos y desean controlar su gasto. Según datos
de la AEVI (Asociación Española de Videojuegos):

\begin{itemize}
\tightlist
\item
  España tiene más de \textbf{16 millones de jugadores} de videojuegos
\item
  El gasto medio anual en videojuegos es de \textbf{150€ por persona}
\item
  El \textbf{68\% de los gamers} tienen entre 16 y 45 años
\end{itemize}

Esto nos da una población objetivo estimada de aproximadamente
\textbf{10,8 millones de personas} en España.

\hypertarget{cuxe1lculo-de-la-muestra-poblacional-necesaria-para-la-encuesta}{%
\subsection{4.4.- Cálculo de la muestra poblacional necesaria para la
encuesta}\label{cuxe1lculo-de-la-muestra-poblacional-necesaria-para-la-encuesta}}

Para una población de 10,8 millones, con un nivel de confianza del 95\%,
un margen de error del 5\% y una proporción esperada del 50\% (máxima
varianza):

\begin{verbatim}
n = (Z² × p × q) / E²
n = (1,96² × 0,5 × 0,5) / 0,05²
n = 384,16
\end{verbatim}

Se necesitarían \textbf{385 encuestas} para obtener resultados
estadísticamente significativos. Para este proyecto, se realizó una
encuesta piloto con una muestra de \textbf{40 personas}.

\hypertarget{encuesta-y-resultados}{%
\subsection{4.5.- Encuesta y resultados}\label{encuesta-y-resultados}}

\textbf{Pregunta 1: ¿Cuánto gastas aproximadamente al mes en
videojuegos?}

\begin{longtable}[]{@{}ll@{}}
\toprule\noalign{}
Rango & Porcentaje \\
\midrule\noalign{}
\endhead
\bottomrule\noalign{}
\endlastfoot
0-20€ & 35\% \\
21-50€ & 40\% \\
51-100€ & 15\% \\
Más de 100€ & 10\% \\
\end{longtable}

\textbf{Pregunta 2: ¿Llevas algún control de tus gastos en videojuegos?}

\begin{longtable}[]{@{}ll@{}}
\toprule\noalign{}
Respuesta & Porcentaje \\
\midrule\noalign{}
\endhead
\bottomrule\noalign{}
\endlastfoot
Sí, de forma manual & 15\% \\
Sí, con alguna app & 10\% \\
No, pero me gustaría & 55\% \\
No, no me interesa & 20\% \\
\end{longtable}

\textbf{Pregunta 3: ¿Te interesaría una app que te ayude a controlar el
presupuesto de gaming?}

\begin{longtable}[]{@{}ll@{}}
\toprule\noalign{}
Respuesta & Porcentaje \\
\midrule\noalign{}
\endhead
\bottomrule\noalign{}
\endlastfoot
Sí, definitivamente & 60\% \\
Tal vez & 25\% \\
No me interesa & 15\% \\
\end{longtable}

\textbf{Pregunta 4: ¿Qué funcionalidades te resultarían más útiles?}

\begin{longtable}[]{@{}ll@{}}
\toprule\noalign{}
Funcionalidad & Porcentaje \\
\midrule\noalign{}
\endhead
\bottomrule\noalign{}
\endlastfoot
Control de gastos mensual & 85\% \\
Lista de deseos con precios & 70\% \\
Catálogo de juegos & 55\% \\
Recomendaciones financieras & 45\% \\
Alertas de ofertas & 75\% \\
Gráficos de tendencia & 50\% \\
\end{longtable}

\textbf{Pregunta 5: ¿Usas alguna app de finanzas personales
actualmente?}

\begin{longtable}[]{@{}ll@{}}
\toprule\noalign{}
Respuesta & Porcentaje \\
\midrule\noalign{}
\endhead
\bottomrule\noalign{}
\endlastfoot
Sí & 20\% \\
No, pero me gustaría & 35\% \\
No, no me interesa & 45\% \\
\end{longtable}

\textbf{Conclusiones de la encuesta}: - El \textbf{75\% de los
encuestados} gasta entre 0 y 50€ mensuales en videojuegos - El
\textbf{65\% no lleva ningún control} de sus gastos en gaming - El
\textbf{85\% estaría interesado} en una app de control de presupuesto
gaming - Las funcionalidades más demandadas son el \textbf{control de
gastos} (85\%) y las \textbf{alertas de ofertas} (75\%)

\hypertarget{planificaciuxf3n-del-tiempo}{%
\subsection{4.6.- Planificación del
tiempo}\label{planificaciuxf3n-del-tiempo}}

El desarrollo de AgenteGamer se planificó siguiendo una metodología ágil
adaptada al contexto de un proyecto de fin de grado. A continuación se
presenta la planificación temporal estimada para las diferentes fases
del proyecto.

\textbf{Fase 1 - Planificación y diseño (Semanas 1-2):}

\begin{itemize}
\tightlist
\item
  Análisis de requisitos y especificaciones
\item
  Diseño de arquitectura y estructura de datos
\item
  Diseño de interfaces de usuario (wireframes)
\item
  Estudio de APIs disponibles (RAWG)
\item
  Planificación de pruebas
\end{itemize}

\textbf{Fase 2 - Desarrollo básico (Semanas 2-10):}

\begin{itemize}
\tightlist
\item
  Implementación de autenticación con Firebase
\item
  Creación de base de datos local con Room
\item
  Desarrollo del módulo de gestión de gastos
\item
  Implementación del catálogo de juegos básico
\end{itemize}

\textbf{Fase 3 - Desarrollo avanzado (Semanas 8-16):}

\begin{itemize}
\tightlist
\item
  Implementación del sistema de recomendaciones del Agente
\item
  Desarrollo de la wishlist con evaluación financiera
\item
  Implementación de seguimiento de lanzamientos próximos
\item
  Integración con Firebase Firestore para sincronización
\end{itemize}

\textbf{Fase 4 - Pruebas y refinamiento (Semanas 16-19):}

\begin{itemize}
\tightlist
\item
  Pruebas unitarias y de integración
\item
  Pruebas de usuario y validación
\item
  Corrección de errores
\item
  Optimización de rendimiento
\end{itemize}

\textbf{Fase 5 - Documentación y presentación (Semanas 18-20):}

\begin{itemize}
\tightlist
\item
  Redacción de documentación técnica
\item
  Elaboración del manual de usuario
\item
  Preparación de materiales para la presentación
\item
  Revisión final y corrección
\end{itemize}

\newpage

\begin{center}
\includegraphics[width=\textwidth,height=\textheight,keepaspectratio]{/home/javier/proyectos/AgenteGamer/AgenteGamer_APP/docs/img/gantt-planificacion.png}
\end{center}

\newpage

\hypertarget{estimaciuxf3n-de-costes}{%
\subsection{4.7.- Estimación de costes}\label{estimaciuxf3n-de-costes}}

La estimación de costes del proyecto se divide en costes de desarrollo,
costes de infraestructura y costes operativos recurrentes.

\textbf{Costes de desarrollo:}

El coste principal es el tiempo de desarrollo del estudiante.
Considerando una estimación de 400 horas de trabajo efectivo a un coste
hora equivalente de 15 euros (basado en el salario mínimo profesional
para desarrolladores junior), el coste de desarrollo se estima en
\textbf{6.000 euros}.

\textbf{Costes de infraestructura:}

\begin{longtable}[]{@{}lll@{}}
\toprule\noalign{}
Concepto & Coste mensual & Coste anual \\
\midrule\noalign{}
\endhead
\bottomrule\noalign{}
\endlastfoot
Firebase (Spark Plan gratuito) & 0 € & 0 € \\
Dominio y hosting web (opcional) & 5 € & 60 € \\
API RAWG (plan gratuito hasta 20.000 req/mes) & 0 € & 0 € \\
\textbf{Total infraestructura} & \textbf{5 €} & \textbf{60 €} \\
\end{longtable}

\textbf{Costes operativos (año 1):}

\begin{longtable}[]{@{}ll@{}}
\toprule\noalign{}
Concepto & Coste \\
\midrule\noalign{}
\endhead
\bottomrule\noalign{}
\endlastfoot
Google Play Developer (una vez) & 25 € \\
Licencia de software (Android Studio - gratuito) & 0 € \\
Mantenimiento y actualizaciones & 200 € \\
\textbf{Total primer año} & \textbf{225 €} \\
\end{longtable}

\textbf{Coste total estimado del proyecto:} 6.285 euros (considerando el
primer año completo).

\hypertarget{estimaciuxf3n-de-ingresos}{%
\subsection{4.8.- Estimación de
ingresos}\label{estimaciuxf3n-de-ingresos}}

La estimación de ingresos se basa en las proyecciones de conversión de
usuarios gratuitos a premium y en el número total de usuarios esperados.

\textbf{Escenario conservador:}

\begin{itemize}
\tightlist
\item
  Usuarios totales esperados (año 1): 5.000 descargas
\item
  Tasa de conversión a premium: 3\%
\item
  Precio medio mensual: 2,99 €
\item
  Ingresos mensuales estimados: 5.000 × 0,03 × 2,99 = \textbf{448,50 €}
\item
  Ingresos anuales estimados: 5.382 €
\end{itemize}

\textbf{Escenario optimista:}

\begin{itemize}
\tightlist
\item
  Usuarios totales esperados (año 1): 20.000 descargas
\item
  Tasa de conversión a premium: 5\%
\item
  Precio medio mensual: 2,99 €
\item
  Ingresos mensuales estimados: 20.000 × 0,05 × 2,99 = \textbf{2.990 €}
\item
  Ingresos anuales estimados: 35.880 €
\end{itemize}

\textbf{Escenario pesimista:}

\begin{itemize}
\tightlist
\item
  Usuarios totales esperados (año 1): 1.000 descargas
\item
  Tasa de conversión a premium: 2\%
\item
  Precio medio mensual: 2,99 €
\item
  Ingresos mensuales estimados: 1.000 × 0,02 × 2,99 = \textbf{59,80 €}
\item
  Ingresos anuales estimados: 717,60 €
\end{itemize}

\vspace{1cm}

\begin{center}
\includegraphics[width=\textwidth]{/home/javier/proyectos/AgenteGamer/AgenteGamer_APP/docs/img/proyeccion-ingresos.png}
\end{center}

\newpage

\hypertarget{conclusiones}{%
\subsection{4.9.- Conclusiones}\label{conclusiones}}

El estudio de viabilidad demuestra que AgenteGamer representa un
proyecto viable tanto técnica como económicamente. La necesidad de
mercado está respaldada por la ausencia de herramientas integradoras que
combinen gestión financiera con catálogo de juegos y recomendaciones
personalizadas. Los resultados de la encuesta confirman que la mayoría
de los gamers no llevan un control de sus gastos y estarían interesados
en una solución específica.

\textbf{Viabilidad técnica:} El proyecto es realizable con las
tecnologías seleccionadas (Java, Android SDK, Room, Firebase, Retrofit),
todas ellas bien documentadas y con comunidades activas de soporte. La
arquitectura Clean Architecture + MVVM permite un desarrollo ordenado y
mantenible.

\textbf{Viabilidad económica:} El modelo de negocio basado en
suscripción ofrece un potencial de ingresos interesante incluso en
escenarios conservadores. Los costes de infraestructura son mínimos
gracias a los planes gratuitos de Firebase y RAWG, lo que permite
mantener la rentabilidad desde los primeros usuarios premium. El punto
de equilibrio es alcanzable con una base de usuarios moderada.

\textbf{Análisis de riesgos:} Se identifican las principales
incertidumbres: la dependencia de una API externa (RAWG), la competencia
con herramientas establecidas de seguimiento de precios, y la necesidad
de construir confianza con los usuarios respecto al manejo de datos
financieros. Estos riesgos se mitigan mediante la implementación de
funcionalidades offline, una propuesta de valor diferenciada, y
políticas de privacidad transparentes.

\newpage

\hypertarget{arquitectura-de-la-aplicaciuxf3n}{%
\section{5.- Arquitectura de la
Aplicación}\label{arquitectura-de-la-aplicaciuxf3n}}

\hypertarget{recursos-y-tecnologuxedas-empleadas}{%
\subsection{5.1.- Recursos y tecnologías
empleadas}\label{recursos-y-tecnologuxedas-empleadas}}

\begin{longtable}[]{@{}
  >{\raggedright\arraybackslash}p{(\columnwidth - 4\tabcolsep) * \real{0.3438}}
  >{\raggedright\arraybackslash}p{(\columnwidth - 4\tabcolsep) * \real{0.3750}}
  >{\raggedright\arraybackslash}p{(\columnwidth - 4\tabcolsep) * \real{0.2812}}@{}}
\toprule\noalign{}
\begin{minipage}[b]{\linewidth}\raggedright
Categoría
\end{minipage} & \begin{minipage}[b]{\linewidth}\raggedright
Tecnología
\end{minipage} & \begin{minipage}[b]{\linewidth}\raggedright
Versión
\end{minipage} \\
\midrule\noalign{}
\endhead
\bottomrule\noalign{}
\endlastfoot
\textbf{Lenguaje} & Java & 17 \\
\textbf{Build System} & Gradle (Kotlin DSL) & - \\
\textbf{SDK} & minSdk 23, compileSdk 36, targetSdk 36 & - \\
\textbf{Arquitectura} & Clean Architecture + MVVM & - \\
\textbf{Inyección de dependencias} & Dagger Hilt & 2.51.1 \\
\textbf{Base de datos local} & Room & 2.8.4 \\
\textbf{Base de datos cloud} & Firestore & Firebase BOM 34.7.0 \\
\textbf{Autenticación} & Firebase Auth & Firebase BOM 34.7.0 \\
\textbf{Networking} & Retrofit + Gson & 2.11.0 \\
\textbf{API externa} & RAWG.io (catálogo de juegos) & - \\
\textbf{Logging HTTP} & OkHttp Logging Interceptor & 4.12.0 \\
\textbf{Gráficos} & MPAndroidChart (PieChart, LineChart) & 3.1.0 \\
\textbf{Carga de imágenes} & Glide & 4.16.0 \\
\textbf{Background tasks} & WorkManager + Hilt Worker & 2.11.0 \\
\textbf{UI Components} & Material Components, BottomNavigationView,
DrawerLayout & 1.13.0 \\
\textbf{Splash Screen} & Splash Screen API & 1.0.1 \\
\textbf{Testing} & JUnit, Mockito, Espresso & 4.13.2 / 5.11.0 / 3.7.0 \\
\end{longtable}

\newpage

\begin{center}
\includegraphics[width=\textwidth]{/home/javier/proyectos/AgenteGamer/AgenteGamer_APP/docs/img/stack-tecnologico.png}
\end{center}

\newpage

\hypertarget{arquitectura-tuxe9cnica}{%
\subsection{5.2.- Arquitectura técnica}\label{arquitectura-tuxe9cnica}}

La aplicación sigue el patrón \textbf{Clean Architecture} con
\textbf{MVVM} (Model-View-ViewModel), utilizando \textbf{Dagger Hilt}
para la inyección de dependencias, \textbf{LiveData} para la reactividad
de la UI y \textbf{Room} como base de datos local con scoping por
userId.

La sincronización con la nube se realiza mediante \textbf{Firebase Auth}
para autenticación y \textbf{Firestore} para datos de usuario
(presupuesto, moneda, perfil). Los gastos, wishlist y lanzamientos se
almacenan localmente en \textbf{Room} con acceso offline.

La integración con la API externa de \textbf{RAWG} se realiza mediante
\textbf{Retrofit + Gson}, con precarga de datos durante la splash screen
para mejorar el rendimiento percibido.

\textbf{Estructura del proyecto:}

\begin{verbatim}
com.miapp.agentegamer/
├── data/
│   ├── local/          # Room entities, DAOs, database
│   ├── remote/         # Retrofit API, DTOs
│   ├── repository/     # Repository implementations
│   └── worker/         # WorkManager workers
├── domain/
│   ├── model/          # Domain models (Gasto, SistemaFinanciero, etc.)
│   ├── repository/     # Repository interfaces
│   └── usecase/        # Use cases
├── di/                 # Hilt modules
├── ui/
│   ├── splash/         # SplashActivity
│   ├── auth/           # Login/Register
│   ├── main/           # Dashboard
│   ├── gastos/         # Expenses
│   ├── games/          # Game catalog
│   ├── wishlist/       # Wishlist
│   ├── lanzamientos/   # Upcoming releases
│   ├── perfil/         # Profile
│   ├── ajustes/        # Settings
│   ├── adapter/        # RecyclerView adapters
│   ├── viewmodel/      # ViewModels
│   └── common/         # BaseNavActivity, utilities
└── util/               # Helpers (MoneyUtils, PeriodoFinancieroUtils, PlatformUtils)
\end{verbatim}

\vspace*{\fill}

\begin{center}
\includegraphics[width=\textwidth]{/home/javier/proyectos/AgenteGamer/AgenteGamer_APP/docs/img/clean-architecture.png}
\end{center}

\vspace*{\fill}

\newpage

\vspace*{\fill}

\begin{center}
\includegraphics[width=\textwidth,height=\textheight,keepaspectratio]{/home/javier/proyectos/AgenteGamer/AgenteGamer_APP/docs/img/diagrama_clases_rotada.png}
\end{center}

\vspace*{\fill}

\newpage

\begin{center}
\includegraphics[width=\textwidth]{/home/javier/proyectos/AgenteGamer/AgenteGamer_APP/docs/img/diagrama_er.png}
\end{center}

\newpage

\hypertarget{diseuxf1o-inicial}{%
\subsection{5.3.- Diseño inicial}\label{diseuxf1o-inicial}}

El diseño inicial contemplaba las siguientes pantallas:

\textbf{Pantalla de Splash}: Logo de AgenteGamer con tagline ``Tu
asistente gaming financiero''. Precarga de datos de juegos en segundo
plano durante 1.5 segundos.

\textbf{Pantalla de Login/Registro}: Formulario simple con email y
contraseña, autenticación mediante Firebase Auth. Registro incluye
nombre y presupuesto inicial.

\textbf{Dashboard Principal}: - Card grande con total gastado y estado
financiero - Dos cards laterales: presupuesto y restante - Gráfico
circular de gastos por categoría - Lista de gastos recientes

\textbf{Gestión de Gastos}: Lista de gastos con opción de añadir, editar
y eliminar. Filtro por fecha.

\textbf{Catálogo de Juegos}: Grid de 2 columnas con búsqueda por nombre.
Integración con RAWG API.

\textbf{Wishlist}: Lista de juegos deseados con precio estimado y coste
total.

\textbf{Lanzamientos}: Lista de juegos próximos con cuenta regresiva.

\textbf{Configuración}: Selector de moneda y ajuste de presupuesto.

\textbf{Perfil}: Información del usuario con opción de editar nombre.

\begin{center}
\includegraphics[width=\textwidth]{/home/javier/proyectos/AgenteGamer/AgenteGamer_APP/docs/img/Sketch_AgenteGamer.png}
\end{center}

\newpage

\hypertarget{diseuxf1o-final}{%
\subsection{5.4.- Diseño final}\label{diseuxf1o-final}}

El diseño final implementado mantiene la estructura inicial pero con
mejoras significativas:

\textbf{UI Revolut-style}: Se adoptó un diseño dark theme inspirado en
la app Revolut, con cards redondeadas (20dp radius), tipografía Manrope
y acentos en verde. La estética gaming se logra con colores oscuros
(midnight blue, slate gray) y acentos verdes.

\textbf{Bottom Navigation Bar}: Se añadió una barra de navegación
inferior con 5 secciones (Home, Gastos, Juegos, Wishlist, Lanzamientos)
para navegación fluida entre secciones sin necesidad de usar el menú
lateral. Implementada mediante una clase base abstracta
\texttt{BaseNavActivity} que centraliza la lógica de navegación con
\texttt{FLAG\_ACTIVITY\_CLEAR\_TOP\ \textbar{}\ FLAG\_ACTIVITY\_SINGLE\_TOP}
para evitar duplicación de actividades. Incluye corrección de selección
visual en \texttt{onResume()} para que el ítem correcto quede resaltado
al volver de otra actividad.

\textbf{Splash Screen con Preloading}: Se implementó una splash screen
que precarga datos de juegos desde la API de RAWG durante el arranque
(mediante \texttt{AgenteGamerApplication.preloadGames()}), reduciendo el
tiempo de carga percibido cuando el usuario llega al catálogo.

\textbf{Gráfico de tendencias}: Se añadió un LineChart mensual que
muestra la evolución de gastos a lo largo del tiempo, con indicador de
tendencia vs mes anterior (porcentaje con flecha y color según
dirección).

\textbf{Recomendaciones del Agente}: Un card dedicado con
recomendaciones personalizadas del agente financiero basadas en los
patrones de gasto del usuario. El sistema evalúa el estado financiero
(verde/amarillo/rojo) según umbrales del presupuesto (50\% verde, 80\%
amarillo, \textgreater80\% rojo) y genera mensajes contextuales.

\textbf{Moneda reactiva}: El selector de moneda en Ajustes se propaga
reactivamente a todas las pantallas mediante LiveData y snapshot
listeners de Firestore. Soporta EUR, USD y GBP con conversión automática
usando tasas fijas.

\textbf{Evaluación financiera de wishlist}: Cada juego de la wishlist
recibe una evaluación automática (recomendada/ajustada/no recomendada)
basada en el presupuesto y gasto actual del mes. El cálculo se recalcula
en tiempo real cuando cambia el presupuesto, el gasto o la wishlist,
usando \texttt{MediatorLiveData} que combina tres fuentes de datos.

\textbf{Compra inteligente}: Al comprar un juego desde la wishlist, el
sistema decide automáticamente: si el juego es un lanzamiento futuro
(fecha \textgreater{} hoy), se guarda en la tabla de Lanzamientos; si ya
está lanzado, se crea un Gasto y se elimina de la wishlist.

\textbf{Notificaciones en segundo plano}: Dos workers de WorkManager
(\texttt{SistemaFinancieroWorker} y \texttt{LanzamientoWorker}) se
ejecutan periódicamente para validar compras contra el presupuesto y
notificar sobre lanzamientos próximos (0-7 días).

\textbf{Menú lateral (Drawer)}: Complementa la bottom navigation con
acceso a Perfil, Ajustes y Cerrar sesión. Incluye header con nombre,
email, rol y fecha de creación del usuario.

\textbf{Scoping por userId}: Todas las consultas de Room están filtradas
por \texttt{userId} de Firebase, garantizando que cada usuario solo ve
sus propios datos.

\newpage

\begin{center}
\includegraphics[width=\textwidth]{/home/javier/proyectos/AgenteGamer/AgenteGamer_APP/docs/img/wireframes-navegacion.png}
\end{center}

\newpage

\hypertarget{futuras-mejoras-y-conclusiones}{%
\section{6.- Futuras Mejoras y
Conclusiones}\label{futuras-mejoras-y-conclusiones}}

Como nunca estoy realmente satisfecho al 100\% con mi trabajo y siempre
pienso que debe haber algo que pueda mejorar, voy a indicar varias cosas
que me hubiera gustado implementar:

\textbf{Notificaciones push para ofertas}: Integrar un sistema de
alertas que notifique al usuario cuando un juego de su wishlist baje de
precio en tiendas como Steam, Epic Games o GOG. Actualmente las
notificaciones solo alertan sobre lanzamientos próximos y compras no
recomendadas.

\textbf{Sincronización con cuentas bancarias}: Mediante APIs como Plaid
o Tink, permitir la importación automática de transacciones relacionadas
con gaming. Actualmente los gastos se registran manualmente.

\textbf{Modo comunidad}: Permitir compartir presupuestos y logros
financieros con amigos, creando un componente social que incentive el
ahorro.

\textbf{Presupuestos por categoría}: Permitir establecer límites de
gasto por categoría (juegos, suscripciones, microtransacciones) con
alertas cuando se acerque al límite. Actualmente solo hay un presupuesto
mensual global.

\textbf{Integración con calendarios de ofertas}: Sincronizar con eventos
como las rebajas de Steam, Black Friday o las ofertas de PlayStation
Store.

\textbf{Modo offline completo}: Mejorar la experiencia offline para que
todas las funcionalidades estén disponibles sin conexión, sincronizando
cuando se reconecte. Actualmente el catálogo de juegos requiere conexión
a la API de RAWG.

\textbf{Análisis de hábitos de compra}: Usar machine learning para
identificar patrones de compra impulsiva y ofrecer recomendaciones
preventivas.

\textbf{Soporte para múltiples perfiles}: Permitir gestionar
presupuestos de diferentes miembros de la familia desde una misma
cuenta.

\textbf{Migraciones de Room no destructivas}: Actualmente la base de
datos usa \texttt{fallbackToDestructiveMigration()}, lo que significa
que cualquier cambio de schema elimina todos los datos locales.
Implementar migraciones progresivas sería esencial para producción.

\textbf{Tests de integración}: Ampliar la cobertura de tests unitarios
actuales con tests de integración que verifiquen el flujo completo de
datos entre Room, ViewModel y UI.

\newpage

\hypertarget{pruebas-y-resultados}{%
\section{7.- Pruebas y Resultados}\label{pruebas-y-resultados}}

\hypertarget{prueba-1-autenticaciuxf3n---registro}{%
\subsection{Prueba 1: Autenticación -
Registro}\label{prueba-1-autenticaciuxf3n---registro}}

\textbf{Escenario}: Registro de nuevo usuario

\textbf{Pasos}: 1. Abrir la aplicación (Splash → Login) 2. Pulsar
``Registrarse'' 3. Introducir email, contraseña, nombre y presupuesto
mensual 4. Pulsar ``Crear cuenta''

\textbf{Resultado esperado}: El usuario se registra en Firebase Auth, se
crea su perfil en Firestore con nombre, presupuesto, rol ``USER'' y
fecha de creación. Se redirige al Dashboard.

\textbf{Resultado obtenido}: Correcto. El usuario se registra y accede
al Dashboard principal.

\hypertarget{prueba-2-autenticaciuxf3n---login}{%
\subsection{Prueba 2: Autenticación -
Login}\label{prueba-2-autenticaciuxf3n---login}}

\textbf{Escenario}: Inicio de sesión con credenciales válidas e
inválidas

\textbf{Pasos}: 1. Abrir la aplicación 2. Introducir email y contraseña
correctos 3. Pulsar ``Iniciar sesión''

\textbf{Resultado esperado}: Navega a MainActivity. Si las credenciales
son incorrectas, muestra mensaje de error específico (red para errores
de red, credenciales para datos incorrectos).

\textbf{Resultado obtenido}: Correcto. Login exitoso redirige al
Dashboard. Credenciales incorrectas muestran toast apropiado.

\hypertarget{prueba-3-dashboard---resumen-financiero}{%
\subsection{Prueba 3: Dashboard - Resumen
financiero}\label{prueba-3-dashboard---resumen-financiero}}

\textbf{Escenario}: Visualización del estado financiero

\textbf{Pasos}: 1. Iniciar sesión como usuario con gastos registrados 2.
Observar el Dashboard

\textbf{Resultado esperado}: Se muestra total gastado, presupuesto,
restante, gráfico circular de categorías, gráfico de tendencia mensual,
recomendaciones del agente y lista de últimos gastos.

\textbf{Resultado obtenido}: Correcto. Todos los componentes del
dashboard se cargan y actualizan reactivamente.

\hypertarget{prueba-4-gestiuxf3n-de-gastos}{%
\subsection{Prueba 4: Gestión de
gastos}\label{prueba-4-gestiuxf3n-de-gastos}}

\textbf{Escenario}: Añadir un nuevo gasto

\textbf{Pasos}: 1. Desde el Dashboard, pulsar ``Gastos'' en la barra
inferior 2. (Solo admin) Pulsar FAB para agregar gasto de prueba

\textbf{Resultado esperado}: El gasto se guarda en Room y se actualiza
la lista con el nuevo gasto. El total se recalcula.

\textbf{Resultado obtenido}: Correcto. El gasto se guarda y la lista se
actualiza automáticamente.

\hypertarget{prueba-5-catuxe1logo-de-juegos}{%
\subsection{Prueba 5: Catálogo de
juegos}\label{prueba-5-catuxe1logo-de-juegos}}

\textbf{Escenario}: Buscar un juego en el catálogo

\textbf{Pasos}: 1. Ir a la sección Juegos 2. Introducir ``Zelda'' en la
barra de búsqueda 3. Esperar los resultados

\textbf{Resultado esperado}: Se muestran juegos con ``Zelda'' en el
nombre, obtenidos de la API RAWG. La búsqueda tiene debounce de 500ms.

\textbf{Resultado obtenido}: Correcto. Los resultados se muestran en un
grid de 2 columnas con paginación infinita.

\hypertarget{prueba-6-wishlist---ediciuxf3n-de-precio}{%
\subsection{Prueba 6: Wishlist - Edición de
precio}\label{prueba-6-wishlist---ediciuxf3n-de-precio}}

\textbf{Escenario}: Editar el precio estimado de un juego en la wishlist

\textbf{Pasos}: 1. Ir a la sección Wishlist 2. Pulsar sobre un juego 3.
Pulsar ``Editar precio'' 4. Introducir nuevo precio y guardar 5. Cerrar
y volver a abrir el juego

\textbf{Resultado esperado}: El nuevo precio se muestra correctamente al
reabrir el detalle.

\textbf{Resultado obtenido}: Correcto. El precio se persiste en Room y
se muestra actualizado.

\hypertarget{prueba-7-navegaciuxf3n-inferior}{%
\subsection{Prueba 7: Navegación
inferior}\label{prueba-7-navegaciuxf3n-inferior}}

\textbf{Escenario}: Navegar entre secciones usando la Bottom Navigation

\textbf{Pasos}: 1. Estar en el Dashboard (Home) 2. Pulsar ``Juegos'' en
la barra inferior 3. Pulsar ``Home'' para volver

\textbf{Resultado esperado}: La navegación es fluida y el ítem correcto
queda resaltado en verde.

\textbf{Resultado obtenido}: Correcto. La selección del bottom nav se
resetea correctamente al volver gracias al \texttt{onResume()} en
\texttt{BaseNavActivity}.

\hypertarget{prueba-8-cambio-de-moneda}{%
\subsection{Prueba 8: Cambio de
moneda}\label{prueba-8-cambio-de-moneda}}

\textbf{Escenario}: Cambiar la moneda en Ajustes

\textbf{Pasos}: 1. Ir a Ajustes desde el menú lateral 2. Cambiar moneda
de EUR a USD 3. Volver al Dashboard

\textbf{Resultado esperado}: Todos los importes se muestran en USD con
la conversión aplicada (tasa 1.08).

\textbf{Resultado obtenido}: Correcto. La moneda se propaga
reactivamente a todas las pantallas.

\hypertarget{prueba-9-perfil---ediciuxf3n-de-nombre}{%
\subsection{Prueba 9: Perfil - Edición de
nombre}\label{prueba-9-perfil---ediciuxf3n-de-nombre}}

\textbf{Escenario}: Editar el nombre del usuario

\textbf{Pasos}: 1. Ir a Perfil desde el menú lateral 2. Pulsar el botón
de editar nombre 3. Introducir nuevo nombre y guardar

\textbf{Resultado esperado}: El nombre se actualiza en Firestore y se
refleja en el perfil.

\textbf{Resultado obtenido}: Correcto. El nombre se actualiza con
indicador de progreso durante la operación.

\hypertarget{prueba-10-ajustes---impacto-en-wishlist}{%
\subsection{Prueba 10: Ajustes - Impacto en
wishlist}\label{prueba-10-ajustes---impacto-en-wishlist}}

\textbf{Escenario}: Cambiar presupuesto y ver impacto en wishlist

\textbf{Pasos}: 1. Ir a Ajustes 2. Modificar el presupuesto mensual 3.
Observar los contadores de impacto (recomendados, ajustados, no
recomendados)

\textbf{Resultado esperado}: Los contadores se recalculan en tiempo real
según el nuevo presupuesto.

\textbf{Resultado obtenido}: Correcto. El impacto se recalcula con cada
cambio en el campo de presupuesto.

\newpage

\hypertarget{tests-unitarios-junit-mockito}{%
\subsection{7.2.- Tests Unitarios (JUnit +
Mockito)}\label{tests-unitarios-junit-mockito}}

La aplicación cuenta con \textbf{10 tests unitarios} que verifican la
lógica de negocio de forma aislada, sin necesidad de un dispositivo
Android. Se ejecutan en la JVM local usando JUnit 4 y Mockito para
mocking de dependencias.

\textbf{Herramientas de testing empleadas:}

\begin{longtable}[]{@{}
  >{\raggedright\arraybackslash}p{(\columnwidth - 4\tabcolsep) * \real{0.3939}}
  >{\raggedright\arraybackslash}p{(\columnwidth - 4\tabcolsep) * \real{0.2727}}
  >{\raggedright\arraybackslash}p{(\columnwidth - 4\tabcolsep) * \real{0.3333}}@{}}
\toprule\noalign{}
\begin{minipage}[b]{\linewidth}\raggedright
Herramienta
\end{minipage} & \begin{minipage}[b]{\linewidth}\raggedright
Versión
\end{minipage} & \begin{minipage}[b]{\linewidth}\raggedright
Propósito
\end{minipage} \\
\midrule\noalign{}
\endhead
\bottomrule\noalign{}
\endlastfoot
JUnit 4 & 4.13.2 & Framework de tests unitarios \\
Mockito & 5.11.0 & Mocking de dependencias (repositorios, DAOs) \\
InstantTaskExecutorRule & 2.2.0 & Ejecución síncrona de LiveData en
tests \\
Espresso & 3.7.0 & Tests de instrumentación en dispositivo \\
\end{longtable}

\hypertarget{tests-de-viewmodels}{%
\subsubsection{Tests de ViewModels}\label{tests-de-viewmodels}}

\textbf{GastoViewModelTest} (10 tests): - Verifica que los LiveData no
sean null (\texttt{getListaGastos}, \texttt{getEstadoUI}) - Testea
\texttt{setSistemaFinanciero()} con gastos y con lista vacía - Verifica
que los métodos CRUD (\texttt{insertar}, \texttt{actualizar},
\texttt{borrar}) deleguen correctamente al repositorio mockeado - Tests
de edge cases con valores null - Verifica que \texttt{onCleared()} no
lance excepciones

\textbf{GamesViewModelTest} (11 tests): - Verifica inicialización del
observer de carga en el constructor - Testea
\texttt{cargarJuegosIniciales()} y \texttt{buscarJuegosPaginados()} con
reset y sin reset - Tests de \texttt{obtenerFamiliasPlataformas()} con
null, PlayStation, múltiples plataformas - Tests de
\texttt{getFamilyIcon()} para familias conocidas y desconocidas -
Verifica limpieza de observers en \texttt{onCleared()}

\hypertarget{tests-de-repositorios}{%
\subsubsection{Tests de Repositorios}\label{tests-de-repositorios}}

\textbf{GastoRepositoryImplTest} (8 tests): - Verifica que
\texttt{obtenerGastos()} devuelva el LiveData del DAO - Testea que las
operaciones CRUD se ejecuten en background mediante
\texttt{ExecutorService} (capturando el Runnable con
\texttt{ArgumentCaptor}) - Verifica \texttt{getGastoMesActual()} y
\texttt{getTotalGastadoMesSync()} con callbacks

\hypertarget{tests-de-entidades-room}{%
\subsubsection{Tests de Entidades Room}\label{tests-de-entidades-room}}

\textbf{GastoEntityTest} (7 tests): - Verifica que el constructor
calcule correctamente \texttt{mes} y \texttt{anio} desde la fecha
(enero, diciembre) - Testea que \texttt{setFecha()} actualice mes y anio
automáticamente - Verifica que se preserven nombre y precio al cambiar
la fecha - Test de cruce de año (31 dic → 1 ene) - Tests de getters y
setters

\textbf{WishlistEntityTest} (8 tests): - Verifica asignación correcta de
campos en el constructor - \textbf{Test específico de bug fix}: confirma
que \texttt{precioEstimado} y \texttt{plataforma} no estén
intercambiados (bug corregido previamente) - Tests con precio cero,
strings vacíos, y verificación de Serializable

\hypertarget{tests-de-utilidades}{%
\subsubsection{Tests de Utilidades}\label{tests-de-utilidades}}

\textbf{MoneyUtilsTest} (19 tests): - Formateo con moneda por defecto
(EUR), explícita (EUR, USD, GBP) - Manejo de moneda null, vacía e
inválida - Formateo de valores negativos - Conversión EUR → USD (100 ×
1.08 = 108), EUR → GBP (100 × 0.86 = 86) - Verificación de que format()
convierte Y formatea correctamente (no muestra el monto original)

\textbf{PresupuestoCalculatorTest} (10 tests): - Cálculo de presupuesto
restante: positivo, cero, negativo, sin gastos - Cálculo de porcentaje
gastado (159.97 / 150 = 106.64\%) - Verificación de sobre-presupuesto y
bajo-presupuesto - \textbf{Caso real del usuario}: presupuesto 150€,
gastado 159.97€, restante -9.97€

\textbf{PeriodoFinancieroUtilsTest} (6 tests): - Verifica que
\texttt{getMesActual()} retorne valor entre 1 y 12 - Verifica
coincidencia con \texttt{Calendar.getInstance()} - Verifica que
\texttt{getAnioActual()} sea razonable (\(\geq\) 2024) - Test de
\texttt{inicioMesActual()}: día 1, hora 0:0:0, menor a hora actual

\textbf{PlatformUtilsTest} (13 tests): - Agrupación de plataformas por
familia: PlayStation (PS4, PS5 → ``playstation''), Xbox, Nintendo, PC -
Test con múltiples plataformas simultáneas (4 familias) - Tests con game
null y platforms null - Tests de \texttt{getFamilyIcon()} para cada
familia y casos desconocidos

\hypertarget{tests-de-instrumentaciuxf3n}{%
\subsubsection{Tests de
Instrumentación}\label{tests-de-instrumentaciuxf3n}}

\textbf{ExampleInstrumentedTest} (1 test): - Verifica que el contexto de
la aplicación sea correcto (\texttt{com.miapp.agentegamer}) - Se ejecuta
en dispositivo Android real o emulador

\hypertarget{resumen-de-cobertura}{%
\subsubsection{Resumen de cobertura}\label{resumen-de-cobertura}}

\begin{longtable}[]{@{}
  >{\raggedright\arraybackslash}p{(\columnwidth - 4\tabcolsep) * \real{0.3793}}
  >{\raggedright\arraybackslash}p{(\columnwidth - 4\tabcolsep) * \real{0.2414}}
  >{\raggedright\arraybackslash}p{(\columnwidth - 4\tabcolsep) * \real{0.3793}}@{}}
\toprule\noalign{}
\begin{minipage}[b]{\linewidth}\raggedright
Categoría
\end{minipage} & \begin{minipage}[b]{\linewidth}\raggedright
Tests
\end{minipage} & \begin{minipage}[b]{\linewidth}\raggedright
Qué cubre
\end{minipage} \\
\midrule\noalign{}
\endhead
\bottomrule\noalign{}
\endlastfoot
ViewModels & 21 & GastoViewModel, GamesViewModel \\
Repositorios & 8 & GastoRepositoryImpl \\
Entidades Room & 15 & GastoEntity, WishlistEntity \\
Utilidades & 48 & MoneyUtils, PresupuestoCalculator,
PeriodoFinancieroUtils, PlatformUtils \\
Instrumentación & 1 & Contexto de la app \\
\textbf{Total} & \textbf{93} & \\
\end{longtable}

\newpage

\hypertarget{manual-usuario}{%
\section{8.- Manual Usuario}\label{manual-usuario}}

\hypertarget{primeros-pasos}{%
\subsection{8.1.- Primeros pasos}\label{primeros-pasos}}

\textbf{Registro}: 1. Abre la aplicación 2. Si no tienes cuenta, pulsa
en ``Registrarse'' 3. Introduce tu email, contraseña, nombre y
presupuesto mensual 4. Pulsa ``Crear cuenta'' 5. Serás redirigido al
Dashboard principal

\textbf{Inicio de sesión}: 1. Abre la aplicación 2. Introduce tu email y
contraseña 3. Pulsa ``Iniciar sesión''

\vspace{0.5cm}

\begin{minipage}{0.48\textwidth}
\centering
\includegraphics[width=\textwidth]{/home/javier/proyectos/AgenteGamer/AgenteGamer_APP/docs/img/registro.png}
\end{minipage}
\hfill
\begin{minipage}{0.48\textwidth}
\centering
\includegraphics[width=\textwidth]{/home/javier/proyectos/AgenteGamer/AgenteGamer_APP/docs/img/inicio_de_sesion.png}
\end{minipage}

\vspace{0.5cm}

\newpage

\hypertarget{dashboard-principal}{%
\subsection{8.2.- Dashboard Principal}\label{dashboard-principal}}

El Dashboard es tu centro de control financiero. Aquí encontrarás:

\begin{itemize}
\tightlist
\item
  \textbf{Total gastado}: El importe total que has gastado este mes en
  videojuegos
\item
  \textbf{Presupuesto}: El límite mensual que has establecido
\item
  \textbf{Restante}: Cuánto te queda por gastar este mes
\item
  \textbf{Indicador de estado}: Punto de color (verde/amarillo/rojo) que
  indica tu situación financiera
\item
  \textbf{Tendencia de gastos}: Gráfico de línea que muestra tu
  evolución mensual
\item
  \textbf{Indicador de tendencia}: Flecha con porcentaje comparando con
  el mes anterior
\item
  \textbf{Tu Agente dice}: Recomendaciones personalizadas basadas en tus
  hábitos
\item
  \textbf{Últimos gastos}: Lista horizontal de tus compras más recientes
\item
  \textbf{Gastos por categoría}: Gráfico circular con la distribución de
  tus gastos
\end{itemize}

\vspace{0.5cm}

\begin{minipage}{0.48\textwidth}
\centering
\includegraphics[width=\textwidth]{/home/javier/proyectos/AgenteGamer/AgenteGamer_APP/docs/img/dashboard.png}
\end{minipage}
\hfill
\begin{minipage}{0.48\textwidth}
\centering
\includegraphics[width=\textwidth]{/home/javier/proyectos/AgenteGamer/AgenteGamer_APP/docs/img/dashboard_dos.png}
\end{minipage}

\vspace{0.5cm}

\newpage

\hypertarget{gestiuxf3n-de-gastos}{%
\subsection{8.3.- Gestión de Gastos}\label{gestiuxf3n-de-gastos}}

\textbf{Ver tus gastos}: 1. Pulsa ``Gastos'' en la barra inferior 2.
Verás la lista de todos tus gastos con imagen, nombre, precio y fecha 3.
El total se muestra en la parte superior

\vspace{0.5cm}

\begin{center}

\includegraphics[width=0.6\textwidth,height=0.7\textheight,keepaspectratio]{/home/javier/proyectos/AgenteGamer/AgenteGamer_APP/docs/img/gastos.png}
\end{center}

\vspace{0.5cm}

\newpage

\hypertarget{catuxe1logo-de-juegos}{%
\subsection{8.4.- Catálogo de Juegos}\label{catuxe1logo-de-juegos}}

\textbf{Buscar juegos}: 1. Pulsa ``Juegos'' en la barra inferior 2. Usa
la barra de búsqueda para encontrar un juego (mínimo 3 caracteres) 3.
Los resultados se muestran en un grid de 2 columnas 4. Haz scroll hacia
abajo para cargar más resultados (paginación infinita)

\textbf{Añadir a la wishlist}: 1. Pulsa sobre cualquier juego del
catálogo 2. El juego se añade automáticamente a tu wishlist con un
precio estimado basado en su rating

\vspace{0.5cm}

\begin{center}
\includegraphics[width=0.6\textwidth,height=0.7\textheight,keepaspectratio]{/home/javier/proyectos/AgenteGamer/AgenteGamer_APP/docs/img/catalogo.png}
\end{center}

\vspace{0.5cm}

\newpage

\hypertarget{wishlist}{%
\subsection{8.5.- Wishlist}\label{wishlist}}

\textbf{Ver tu lista de deseos}: 1. Pulsa ``Wishlist'' en la barra
inferior 2. Verás todos los juegos que has añadido con su precio
estimado y evaluación financiera

\textbf{Evaluaciones financieras}: - \textbf{Compra recomendada}: El
precio es bajo comparado con tu presupuesto restante - \textbf{Compra
ajustada}: El precio es moderado, podrías comprarlo pero con precaución
- \textbf{Mejor esperar}: El precio consume gran parte de tu presupuesto
restante - \textbf{No recomendable}: No tienes suficiente presupuesto
para esta compra

\textbf{Editar precio estimado}: 1. Pulsa sobre un juego de la wishlist
2. Pulsa ``Editar precio'' 3. Introduce el nuevo precio y pulsa
``Guardar''

\textbf{Eliminar de la wishlist}: 1. Pulsa sobre el botón de eliminar
junto al juego 2. El juego se elimina de tu lista

\textbf{Comprar un juego}: 1. Pulsa sobre un juego de la wishlist 2.
Pulsa ``Comprar'' 3. Si el juego aún no ha sido lanzado, se moverá a
Lanzamientos 4. Si ya está disponible, se creará un gasto y se eliminará
de la wishlist

\begin{minipage}{0.48\textwidth}
\centering
\includegraphics[width=\textwidth,height=0.6\textheight,keepaspectratio]{/home/javier/proyectos/AgenteGamer/AgenteGamer_APP/docs/img/whislist.png}
\end{minipage}
\hfill
\begin{minipage}{0.48\textwidth}
\centering
\includegraphics[width=\textwidth,height=0.6\textheight,keepaspectratio]{/home/javier/proyectos/AgenteGamer/AgenteGamer_APP/docs/img/wishlist_editor.png}
\end{minipage}

\newpage

\hypertarget{lanzamientos}{%
\subsection{8.6.- Lanzamientos}\label{lanzamientos}}

\textbf{Ver próximos lanzamientos}: 1. Pulsa ``Lanzamientos'' en la
barra inferior 2. Verás una lista de juegos próximos con días restantes,
rating y plataformas

\vspace{0.5cm}

\begin{center}
\includegraphics[width=0.6\textwidth,height=0.7\textheight,keepaspectratio]{/home/javier/proyectos/AgenteGamer/AgenteGamer_APP/docs/img/lanzamientos.png}
\end{center}

\vspace{0.5cm}

\newpage

\hypertarget{perfil}{%
\subsection{8.7.- Perfil}\label{perfil}}

\textbf{Ver tu perfil}: 1. Abre el menú lateral (arriba a la izquierda)
ubicado en la pantalla principal. 2. Pulsa ``Perfil'' 3. Verás tu email,
nombre, presupuesto, fecha de creación y rol

\textbf{Editar tu nombre}: 1. En la pantalla de Perfil, pulsa el botón
de editar junto a tu nombre 2. Introduce el nuevo nombre 3. Pulsa
``Guardar''

\vspace{0.5cm}

\begin{minipage}{0.48\textwidth}
\centering
\includegraphics[width=\textwidth]{/home/javier/proyectos/AgenteGamer/AgenteGamer_APP/docs/img/perfil_uno.png}
\end{minipage}
\hfill
\begin{minipage}{0.48\textwidth}
\centering
\includegraphics[width=\textwidth]{/home/javier/proyectos/AgenteGamer/AgenteGamer_APP/docs/img/perfil_tres.png}
\end{minipage}

\vspace{0.5cm}

\newpage

\hypertarget{configuraciuxf3n}{%
\subsection{8.8.- Configuración}\label{configuraciuxf3n}}

\textbf{Cambiar moneda}: 1. Abre el menú lateral (arriba a la izquierda)
ubicado en la pantalla principal. 2. Ve a ``Ajustes'' 3. Selecciona tu
moneda preferida: EUR, USD o GBP 4. El cambio se aplica automáticamente
a todas las pantallas

\textbf{Ajustar presupuesto}: 1. Ve a ``Ajustes'' 2. Introduce tu nuevo
presupuesto mensual 3. Observa cómo cambia el impacto en tu wishlist en
tiempo real 4. Pulsa ``Guardar''

\vspace{0.5cm}

\begin{center}
\includegraphics[width=0.6\textwidth,height=0.7\textheight,keepaspectratio]{/home/javier/proyectos/AgenteGamer/AgenteGamer_APP/docs/img/perfil_dos.png}
\end{center}

\vspace{0.5cm}

\newpage

\hypertarget{menuxfa-lateral-drawer}{%
\subsection{8.9.- Menú lateral (Drawer)}\label{menuxfa-lateral-drawer}}

El menú lateral se abre deslizando desde el borde izquierdo o pulsando
el icono de hamburguesa en la toolbar. Contiene:

\begin{itemize}
\tightlist
\item
  \textbf{Perfil}: Ver y editar tu información
\item
  \textbf{Juegos}: Ir al catálogo de juegos
\item
  \textbf{Wishlist}: Ir a tu lista de deseos
\item
  \textbf{Gastos}: Ir al historial de gastos
\item
  \textbf{Lanzamientos}: Ver próximos lanzamientos
\item
  \textbf{Ajustes}: Configurar moneda y presupuesto
\item
  \textbf{Cerrar sesión}: Salir de la aplicación
\end{itemize}

\begin{center}
\includegraphics[width=0.6\textwidth,height=0.7\textheight,keepaspectratio]{/home/javier/proyectos/AgenteGamer/AgenteGamer_APP/docs/img/perfil_uno.png}
\end{center}

\newpage

\hypertarget{bibliografuxeda}{%
\section{9.- Bibliografía}\label{bibliografuxeda}}

\begin{itemize}
\tightlist
\item
  \textbf{Documentación de Android Developers.} Android Developers.
  {[}En línea{]}. Disponible en: \url{https://developer.android.com/}
\item
  \textbf{Documentación de Firebase.} Google. {[}En línea{]}. Disponible
  en: \url{https://firebase.google.com/docs}
\item
  \textbf{Documentación de Room Database.} Google. {[}En línea{]}.
  Disponible en:
  \url{https://developer.android.com/jetpack/androidx/releases/room}
\item
  \textbf{Documentación de Dagger Hilt.} Google. {[}En línea{]}.
  Disponible en:
  \url{https://developer.android.com/training/dependency-injection/hilt-android}
\item
  \textbf{RAWG Video Games Database API.} RAWG. {[}En línea{]}.
  Disponible en: \url{https://rawg.io/apidocs}
\item
  \textbf{Documentación de MPAndroidChart.} PhilJay. {[}En línea{]}.
  Disponible en: \url{https://github.com/PhilJay/MPAndroidChart}
\item
  \textbf{Documentación de Retrofit.} Square. {[}En línea{]}. Disponible
  en: \url{https://square.github.io/retrofit/}
\item
  \textbf{Documentación de WorkManager.} Google. {[}En línea{]}.
  Disponible en:
  \url{https://developer.android.com/topic/libraries/architecture/workmanager}
\item
  \textbf{Clean Architecture para aplicaciones Android.} Google
  Codelabs. {[}En línea{]}. Disponible en:
  \url{https://developer.android.com/courses/pathways/android-architecture}
\item
  \textbf{Material Design Guidelines.} Google. {[}En línea{]}.
  Disponible en: \url{https://material.io/design}
\item
  \textbf{Informe anual de la industria del videojuegos 2024.} AEVI -
  Asociación Española de Videojuegos. {[}En línea{]}. Disponible en:
  \url{https://www.aevi.org.es/}
\item
  \textbf{Estadísticas de consumo de videojuegos.} Newzoo. {[}En
  línea{]}. Disponible en: \url{https://newzoo.com/}
\end{itemize}

\newpage

\hypertarget{anexo-i---principales-clases-del-cuxf3digo}{%
\section{10.- Anexo I - Principales clases del
código}\label{anexo-i---principales-clases-del-cuxf3digo}}

\hypertarget{entidades-room}{%
\subsection{Entidades Room}\label{entidades-room}}

\textbf{GastoEntity.java} - Tabla \texttt{gastos}:

\begin{Shaded}
\begin{Highlighting}[]
\AttributeTok{@Entity}\OperatorTok{(}\NormalTok{tableName }\OperatorTok{=} \StringTok{"gastos"}\OperatorTok{,}
\NormalTok{    indices }\OperatorTok{=} \OperatorTok{\{}\AttributeTok{@Index}\OperatorTok{(}\NormalTok{value }\OperatorTok{=} \OperatorTok{\{}\StringTok{"mes"}\OperatorTok{,} \StringTok{"anio"}\OperatorTok{\}),} \AttributeTok{@Index}\OperatorTok{(}\NormalTok{value }\OperatorTok{=} \OperatorTok{\{}\StringTok{"userId"}\OperatorTok{\})\})}
\KeywordTok{public} \KeywordTok{class}\NormalTok{ GastoEntity }\OperatorTok{\{}
    \AttributeTok{@PrimaryKey}\OperatorTok{(}\NormalTok{autoGenerate }\OperatorTok{=} \KeywordTok{true}\OperatorTok{)}
    \KeywordTok{private} \DataTypeTok{int}\NormalTok{ id}\OperatorTok{;}
    \KeywordTok{private} \BuiltInTok{String}\NormalTok{ userId}\OperatorTok{;}
    \KeywordTok{private} \BuiltInTok{String}\NormalTok{ nombreJuego}\OperatorTok{;}
    \KeywordTok{private} \DataTypeTok{double}\NormalTok{ precio}\OperatorTok{;}
    \KeywordTok{private} \DataTypeTok{long}\NormalTok{ fecha}\OperatorTok{;}  \CommentTok{// UNIX epoch en milisegundos}
    \KeywordTok{private} \DataTypeTok{int}\NormalTok{ mes}\OperatorTok{;}     \CommentTok{// Calculado automáticamente}
    \KeywordTok{private} \DataTypeTok{int}\NormalTok{ anio}\OperatorTok{;}    \CommentTok{// Calculado automáticamente}
    \KeywordTok{private} \BuiltInTok{String}\NormalTok{ imagenUrl}\OperatorTok{;}
    \CommentTok{// Getters y setters con recálculo de mes/anio al cambiar fecha}
\OperatorTok{\}}
\end{Highlighting}
\end{Shaded}

\textbf{WishlistEntity.java} - Tabla \texttt{wishlist}:

\begin{Shaded}
\begin{Highlighting}[]
\AttributeTok{@Entity}\OperatorTok{(}\NormalTok{tableName }\OperatorTok{=} \StringTok{"wishlist"}\OperatorTok{)}
\KeywordTok{public} \KeywordTok{class}\NormalTok{ WishlistEntity }\OperatorTok{\{}
    \AttributeTok{@PrimaryKey}\OperatorTok{(}\NormalTok{autoGenerate }\OperatorTok{=} \KeywordTok{true}\OperatorTok{)}
    \KeywordTok{private} \DataTypeTok{int}\NormalTok{ id}\OperatorTok{;}
    \KeywordTok{private} \BuiltInTok{String}\NormalTok{ userId}\OperatorTok{;}
    \KeywordTok{private} \DataTypeTok{int}\NormalTok{ gameId}\OperatorTok{;}
    \KeywordTok{private} \BuiltInTok{String}\NormalTok{ nombre}\OperatorTok{;}
    \KeywordTok{private} \BuiltInTok{String}\NormalTok{ fechaLanzamiento}\OperatorTok{;}
    \KeywordTok{private} \BuiltInTok{String}\NormalTok{ imagenUrl}\OperatorTok{;}
    \KeywordTok{private} \BuiltInTok{String}\NormalTok{ plataforma}\OperatorTok{;}
    \KeywordTok{private} \DataTypeTok{double}\NormalTok{ precioEstimado}\OperatorTok{;}
    \CommentTok{// Getters y setters}
\OperatorTok{\}}
\end{Highlighting}
\end{Shaded}

\textbf{LanzamientoEntity.java} - Tabla \texttt{lanzamientos}:

\begin{Shaded}
\begin{Highlighting}[]
\AttributeTok{@Entity}\OperatorTok{(}\NormalTok{tableName }\OperatorTok{=} \StringTok{"lanzamientos"}\OperatorTok{)}
\KeywordTok{public} \KeywordTok{class}\NormalTok{ LanzamientoEntity }\OperatorTok{\{}
    \AttributeTok{@PrimaryKey}\OperatorTok{(}\NormalTok{autoGenerate }\OperatorTok{=} \KeywordTok{true}\OperatorTok{)}
    \KeywordTok{private} \DataTypeTok{int}\NormalTok{ id}\OperatorTok{;}
    \KeywordTok{private} \BuiltInTok{String}\NormalTok{ userId}\OperatorTok{;}
    \KeywordTok{private} \DataTypeTok{int}\NormalTok{ gameId}\OperatorTok{;}
    \KeywordTok{private} \BuiltInTok{String}\NormalTok{ nombre}\OperatorTok{;}
    \KeywordTok{private} \DataTypeTok{long}\NormalTok{ fechaLanzamiento}\OperatorTok{;}
    \KeywordTok{private} \DataTypeTok{double}\NormalTok{ precio}\OperatorTok{;}
    \KeywordTok{private} \BuiltInTok{String}\NormalTok{ imagenUrl}\OperatorTok{;}
    \KeywordTok{private} \BuiltInTok{String}\NormalTok{ plataformas}\OperatorTok{;}
    \KeywordTok{private} \DataTypeTok{double}\NormalTok{ rating}\OperatorTok{;}
    \CommentTok{// Getters y setters}
\OperatorTok{\}}
\end{Highlighting}
\end{Shaded}

\hypertarget{base-de-datos}{%
\subsection{Base de datos}\label{base-de-datos}}

\textbf{AppDatabase.java}:

\begin{Shaded}
\begin{Highlighting}[]
\AttributeTok{@Database}\OperatorTok{(}\NormalTok{entities }\OperatorTok{=} \OperatorTok{\{}\NormalTok{GastoEntity}\OperatorTok{.}\FunctionTok{class}\OperatorTok{,}\NormalTok{ WishlistEntity}\OperatorTok{.}\FunctionTok{class}\OperatorTok{,}
\NormalTok{    LanzamientoEntity}\OperatorTok{.}\FunctionTok{class}\OperatorTok{\},}\NormalTok{ version }\OperatorTok{=} \DecValTok{9}\OperatorTok{,}\NormalTok{ exportSchema }\OperatorTok{=} \KeywordTok{false}\OperatorTok{)}
\KeywordTok{public} \KeywordTok{abstract} \KeywordTok{class}\NormalTok{ AppDatabase }\KeywordTok{extends}\NormalTok{ RoomDatabase }\OperatorTok{\{}
    \KeywordTok{public} \KeywordTok{abstract}\NormalTok{ GastoDao }\FunctionTok{gastoDao}\OperatorTok{();}
    \KeywordTok{public} \KeywordTok{abstract}\NormalTok{ WishlistDao }\FunctionTok{wishlistDao}\OperatorTok{();}
    \KeywordTok{public} \KeywordTok{abstract}\NormalTok{ LanzamientoDao }\FunctionTok{lanzamientoDao}\OperatorTok{();}
\OperatorTok{\}}
\end{Highlighting}
\end{Shaded}

\hypertarget{daos}{%
\subsection{DAOs}\label{daos}}

\textbf{GastoDao.java} - Operaciones sobre gastos:

\begin{Shaded}
\begin{Highlighting}[]
\AttributeTok{@Dao}
\KeywordTok{public} \KeywordTok{interface}\NormalTok{ GastoDao }\OperatorTok{\{}
    \AttributeTok{@Insert} \DataTypeTok{void} \FunctionTok{insertGasto}\OperatorTok{(}\NormalTok{GastoEntity gasto}\OperatorTok{);}
    \AttributeTok{@Update} \DataTypeTok{void} \FunctionTok{updateGasto}\OperatorTok{(}\NormalTok{GastoEntity gasto}\OperatorTok{);}
    \AttributeTok{@Delete} \DataTypeTok{void} \FunctionTok{deleteGasto}\OperatorTok{(}\NormalTok{GastoEntity gasto}\OperatorTok{);}

    \AttributeTok{@Query}\OperatorTok{(}\StringTok{"SELECT * FROM gastos WHERE userId = :userId ORDER BY fecha DESC"}\OperatorTok{)}
\NormalTok{    LiveData}\OperatorTok{\textless{}}\BuiltInTok{List}\OperatorTok{\textless{}}\NormalTok{GastoEntity}\OperatorTok{\textgreater{}\textgreater{}} \FunctionTok{getAllGastos}\OperatorTok{(}\BuiltInTok{String}\NormalTok{ userId}\OperatorTok{);}

    \AttributeTok{@Query}\OperatorTok{(}\StringTok{"SELECT IFNULL(SUM(precio),0) FROM gastos WHERE userId = :userId "} \OperatorTok{+}
           \StringTok{"AND mes = :mes AND anio = :anio"}\OperatorTok{)}
    \DataTypeTok{double} \FunctionTok{getTotalGastadoMes}\OperatorTok{(}\BuiltInTok{String}\NormalTok{ userId}\OperatorTok{,} \DataTypeTok{int}\NormalTok{ mes}\OperatorTok{,} \DataTypeTok{int}\NormalTok{ anio}\OperatorTok{);}

    \AttributeTok{@Query}\OperatorTok{(}\StringTok{"SELECT strftime(\textquotesingle{}\%Y{-}\%m\textquotesingle{}, fecha/1000, \textquotesingle{}unixepoch\textquotesingle{}) as mes, "} \OperatorTok{+}
           \StringTok{"SUM(precio) as total FROM gastos WHERE userId = :userId "} \OperatorTok{+}
           \StringTok{"AND fecha \textgreater{}= :startDate GROUP BY mes ORDER BY mes ASC LIMIT :limit"}\OperatorTok{)}
    \BuiltInTok{List}\OperatorTok{\textless{}}\NormalTok{MonthlyTotal}\OperatorTok{\textgreater{}} \FunctionTok{getMonthlyTotals}\OperatorTok{(}\BuiltInTok{String}\NormalTok{ userId}\OperatorTok{,} \DataTypeTok{long}\NormalTok{ startDate}\OperatorTok{,} \DataTypeTok{int}\NormalTok{ limit}\OperatorTok{);}

    \AttributeTok{@Query}\OperatorTok{(}\StringTok{"SELECT * FROM gastos WHERE userId = :userId "} \OperatorTok{+}
           \StringTok{"ORDER BY fecha DESC LIMIT :limit"}\OperatorTok{)}
\NormalTok{    LiveData}\OperatorTok{\textless{}}\BuiltInTok{List}\OperatorTok{\textless{}}\NormalTok{GastoEntity}\OperatorTok{\textgreater{}\textgreater{}} \FunctionTok{getRecentGastos}\OperatorTok{(}\BuiltInTok{String}\NormalTok{ userId}\OperatorTok{,} \DataTypeTok{int}\NormalTok{ limit}\OperatorTok{);}
\OperatorTok{\}}
\end{Highlighting}
\end{Shaded}

\textbf{WishlistDao.java} - Operaciones sobre wishlist:

\begin{Shaded}
\begin{Highlighting}[]
\AttributeTok{@Dao}
\KeywordTok{public} \KeywordTok{interface}\NormalTok{ WishlistDao }\OperatorTok{\{}
    \AttributeTok{@Insert}\OperatorTok{(}\NormalTok{onConflict }\OperatorTok{=}\NormalTok{ OnConflictStrategy}\OperatorTok{.}\FunctionTok{REPLACE}\OperatorTok{)}
    \DataTypeTok{void} \FunctionTok{insert}\OperatorTok{(}\NormalTok{WishlistEntity juego}\OperatorTok{);}

    \AttributeTok{@Query}\OperatorTok{(}\StringTok{"SELECT * FROM wishlist WHERE userId = :userId"}\OperatorTok{)}
\NormalTok{    LiveData}\OperatorTok{\textless{}}\BuiltInTok{List}\OperatorTok{\textless{}}\NormalTok{WishlistEntity}\OperatorTok{\textgreater{}\textgreater{}} \FunctionTok{getWishlist}\OperatorTok{(}\BuiltInTok{String}\NormalTok{ userId}\OperatorTok{);}

    \AttributeTok{@Update} \DataTypeTok{void} \FunctionTok{actualizar}\OperatorTok{(}\NormalTok{WishlistEntity juego}\OperatorTok{);}
    \AttributeTok{@Delete} \DataTypeTok{void} \FunctionTok{borrar}\OperatorTok{(}\NormalTok{WishlistEntity juego}\OperatorTok{);}

    \AttributeTok{@Query}\OperatorTok{(}\StringTok{"SELECT * FROM wishlist WHERE userId = :userId"}\OperatorTok{)}
    \BuiltInTok{List}\OperatorTok{\textless{}}\NormalTok{WishlistEntity}\OperatorTok{\textgreater{}} \FunctionTok{getWishlistSync}\OperatorTok{(}\BuiltInTok{String}\NormalTok{ userId}\OperatorTok{);}

    \AttributeTok{@Query}\OperatorTok{(}\StringTok{"SELECT SUM(precioEstimado) FROM wishlist WHERE userId = :userId"}\OperatorTok{)}
    \BuiltInTok{Double} \FunctionTok{getTotalGastado}\OperatorTok{(}\BuiltInTok{String}\NormalTok{ userId}\OperatorTok{);}
\OperatorTok{\}}
\end{Highlighting}
\end{Shaded}

\hypertarget{modelo-de-dominio}{%
\subsection{Modelo de dominio}\label{modelo-de-dominio}}

\textbf{SistemaFinanciero.java} - Lógica de negocio financiera:

\begin{Shaded}
\begin{Highlighting}[]
\KeywordTok{public} \KeywordTok{class}\NormalTok{ SistemaFinanciero }\OperatorTok{\{}
    \CommentTok{// Umbrales: 50\% verde, 80\% amarillo, \textgreater{}80\% rojo}
    \KeywordTok{private} \DataTypeTok{static} \DataTypeTok{final} \DataTypeTok{double}\NormalTok{ PORCENTAJE\_UMBRAL\_VERDE }\OperatorTok{=} \FloatTok{0.5}\OperatorTok{;}
    \KeywordTok{private} \DataTypeTok{static} \DataTypeTok{final} \DataTypeTok{double}\NormalTok{ PORCENTAJE\_UMBRAL\_AMARILLO }\OperatorTok{=} \FloatTok{0.8}\OperatorTok{;}

    \CommentTok{// Umbrales de evaluación de compra}
    \KeywordTok{private} \DataTypeTok{static} \DataTypeTok{final} \DataTypeTok{double}\NormalTok{ PORCENTAJE\_EVALUAR\_RECOMENDADO }\OperatorTok{=} \FloatTok{0.3}\OperatorTok{;}
    \KeywordTok{private} \DataTypeTok{static} \DataTypeTok{final} \DataTypeTok{double}\NormalTok{ PORCENTAJE\_EVALUAR\_AJUSTADO }\OperatorTok{=} \FloatTok{0.6}\OperatorTok{;}

    \CommentTok{// Precios estimados por rating}
    \KeywordTok{private} \DataTypeTok{static} \DataTypeTok{final} \DataTypeTok{double}\NormalTok{ PRECIO\_ALTO }\OperatorTok{=} \FloatTok{69.99}\OperatorTok{;}    \CommentTok{// rating \textgreater{}= 4.5}
    \KeywordTok{private} \DataTypeTok{static} \DataTypeTok{final} \DataTypeTok{double}\NormalTok{ PRECIO\_MEDIO\_ALTO }\OperatorTok{=} \FloatTok{59.99}\OperatorTok{;} \CommentTok{// rating \textgreater{}= 4.0}
    \KeywordTok{private} \DataTypeTok{static} \DataTypeTok{final} \DataTypeTok{double}\NormalTok{ PRECIO\_MEDIO }\OperatorTok{=} \FloatTok{49.99}\OperatorTok{;}   \CommentTok{// rating \textgreater{}= 3.5}
    \KeywordTok{private} \DataTypeTok{static} \DataTypeTok{final} \DataTypeTok{double}\NormalTok{ PRECIO\_BAJO }\OperatorTok{=} \FloatTok{39.99}\OperatorTok{;}    \CommentTok{// rating \textless{} 3.5}

    \KeywordTok{private} \DataTypeTok{final} \DataTypeTok{double}\NormalTok{ presupuestoMensual}\OperatorTok{;}

    \KeywordTok{public}\NormalTok{ EstadoFinanciero }\FunctionTok{obtenerEstado}\OperatorTok{(}\DataTypeTok{double}\NormalTok{ totalGastado}\OperatorTok{)} \OperatorTok{\{}
        \ControlFlowTok{if} \OperatorTok{(}\NormalTok{totalGastado }\OperatorTok{\textless{}}\NormalTok{ presupuestoMensual }\OperatorTok{*}\NormalTok{ PORCENTAJE\_UMBRAL\_VERDE}\OperatorTok{)}
            \ControlFlowTok{return}\NormalTok{ EstadoFinanciero}\OperatorTok{.}\FunctionTok{VERDE}\OperatorTok{;}
        \ControlFlowTok{else} \ControlFlowTok{if} \OperatorTok{(}\NormalTok{totalGastado }\OperatorTok{\textless{}}\NormalTok{ presupuestoMensual }\OperatorTok{*}\NormalTok{ PORCENTAJE\_UMBRAL\_AMARILLO}\OperatorTok{)}
            \ControlFlowTok{return}\NormalTok{ EstadoFinanciero}\OperatorTok{.}\FunctionTok{AMARILLO}\OperatorTok{;}
        \ControlFlowTok{else}
            \ControlFlowTok{return}\NormalTok{ EstadoFinanciero}\OperatorTok{.}\FunctionTok{ROJO}\OperatorTok{;}
    \OperatorTok{\}}

    \KeywordTok{public} \BuiltInTok{String} \FunctionTok{evaluarCompra}\OperatorTok{(}\DataTypeTok{double}\NormalTok{ precioJuego}\OperatorTok{,} \DataTypeTok{double}\NormalTok{ gastoMes}\OperatorTok{)} \OperatorTok{\{}
        \DataTypeTok{double}\NormalTok{ restante }\OperatorTok{=}\NormalTok{ presupuestoMensual }\OperatorTok{{-}}\NormalTok{ gastoMes}\OperatorTok{;}
        \ControlFlowTok{if} \OperatorTok{(}\NormalTok{restante }\OperatorTok{\textless{}} \DecValTok{0}\OperatorTok{)} \ControlFlowTok{return} \StringTok{"No recomendable"}\OperatorTok{;}
        \ControlFlowTok{if} \OperatorTok{(}\NormalTok{precioJuego }\OperatorTok{\textless{}=}\NormalTok{ restante }\OperatorTok{*}\NormalTok{ PORCENTAJE\_EVALUAR\_RECOMENDADO}\OperatorTok{)}
            \ControlFlowTok{return} \StringTok{"Compra recomendada"}\OperatorTok{;}
        \ControlFlowTok{else} \ControlFlowTok{if} \OperatorTok{(}\NormalTok{precioJuego }\OperatorTok{\textless{}=}\NormalTok{ restante }\OperatorTok{*}\NormalTok{ PORCENTAJE\_EVALUAR\_AJUSTADO}\OperatorTok{)}
            \ControlFlowTok{return} \StringTok{"Compra ajustada"}\OperatorTok{;}
        \ControlFlowTok{else} \ControlFlowTok{if} \OperatorTok{(}\NormalTok{precioJuego }\OperatorTok{\textless{}=}\NormalTok{ restante}\OperatorTok{)}
            \ControlFlowTok{return} \StringTok{"Mejor esperar"}\OperatorTok{;}
        \ControlFlowTok{else}
            \ControlFlowTok{return} \StringTok{"No recomendable"}\OperatorTok{;}
    \OperatorTok{\}}

    \KeywordTok{public} \DataTypeTok{double} \FunctionTok{estimarPrecio}\OperatorTok{(}\NormalTok{GameInfo juego}\OperatorTok{)} \OperatorTok{\{}
        \ControlFlowTok{if} \OperatorTok{(}\NormalTok{juego }\OperatorTok{==} \KeywordTok{null} \OperatorTok{||}\NormalTok{ juego}\OperatorTok{.}\FunctionTok{getRating}\OperatorTok{()} \OperatorTok{\textless{}=} \DecValTok{0}\OperatorTok{)} \ControlFlowTok{return}\NormalTok{ PRECIO\_MEDIO}\OperatorTok{;}
        \ControlFlowTok{if} \OperatorTok{(}\NormalTok{juego}\OperatorTok{.}\FunctionTok{getRating}\OperatorTok{()} \OperatorTok{\textgreater{}=} \FloatTok{4.5}\OperatorTok{)} \ControlFlowTok{return}\NormalTok{ PRECIO\_ALTO}\OperatorTok{;}
        \ControlFlowTok{if} \OperatorTok{(}\NormalTok{juego}\OperatorTok{.}\FunctionTok{getRating}\OperatorTok{()} \OperatorTok{\textgreater{}=} \FloatTok{4.0}\OperatorTok{)} \ControlFlowTok{return}\NormalTok{ PRECIO\_MEDIO\_ALTO}\OperatorTok{;}
        \ControlFlowTok{if} \OperatorTok{(}\NormalTok{juego}\OperatorTok{.}\FunctionTok{getRating}\OperatorTok{()} \OperatorTok{\textgreater{}=} \FloatTok{3.5}\OperatorTok{)} \ControlFlowTok{return}\NormalTok{ PRECIO\_MEDIO}\OperatorTok{;}
        \ControlFlowTok{return}\NormalTok{ PRECIO\_BAJO}\OperatorTok{;}
    \OperatorTok{\}}
\OperatorTok{\}}
\end{Highlighting}
\end{Shaded}

\hypertarget{viewmodels}{%
\subsection{ViewModels}\label{viewmodels}}

\textbf{GastoViewModel.java} - Gestión de gastos y dashboard:

\begin{Shaded}
\begin{Highlighting}[]
\AttributeTok{@HiltViewModel}
\KeywordTok{public} \KeywordTok{class}\NormalTok{ GastoViewModel }\KeywordTok{extends}\NormalTok{ AndroidViewModel }\OperatorTok{\{}
    \KeywordTok{private} \DataTypeTok{final}\NormalTok{ GastoRepository repository}\OperatorTok{;}
    \KeywordTok{private} \DataTypeTok{final}\NormalTok{ LiveData}\OperatorTok{\textless{}}\BuiltInTok{List}\OperatorTok{\textless{}}\NormalTok{GastoEntity}\OperatorTok{\textgreater{}\textgreater{}}\NormalTok{ listaGastos}\OperatorTok{;}
    \KeywordTok{private} \DataTypeTok{final}\NormalTok{ MutableLiveData}\OperatorTok{\textless{}}\NormalTok{EstadoFinancieroUI}\OperatorTok{\textgreater{}}\NormalTok{ estadoUI}\OperatorTok{;}
    \KeywordTok{private} \DataTypeTok{final}\NormalTok{ MutableLiveData}\OperatorTok{\textless{}}\BuiltInTok{List}\OperatorTok{\textless{}}\NormalTok{MonthlyExpense}\OperatorTok{\textgreater{}\textgreater{}}\NormalTok{ monthlyExpenses}\OperatorTok{;}
    \KeywordTok{private} \DataTypeTok{final}\NormalTok{ MutableLiveData}\OperatorTok{\textless{}}\NormalTok{TrendResult}\OperatorTok{\textgreater{}}\NormalTok{ trendResult}\OperatorTok{;}
    \KeywordTok{private} \DataTypeTok{final}\NormalTok{ LiveData}\OperatorTok{\textless{}}\BuiltInTok{List}\OperatorTok{\textless{}}\NormalTok{GastoEntity}\OperatorTok{\textgreater{}\textgreater{}}\NormalTok{ recentGastos}\OperatorTok{;}

    \KeywordTok{public} \DataTypeTok{void} \FunctionTok{setSistemaFinanciero}\OperatorTok{(}\NormalTok{SistemaFinanciero sistema}\OperatorTok{)} \OperatorTok{\{}
        \KeywordTok{this}\OperatorTok{.}\FunctionTok{sistemaFinanciero} \OperatorTok{=}\NormalTok{ sistema}\OperatorTok{;}
        \FunctionTok{recalcularEstado}\OperatorTok{(}\NormalTok{listaGastos}\OperatorTok{.}\FunctionTok{getValue}\OperatorTok{());}
        \FunctionTok{loadDashboardData}\OperatorTok{();}
    \OperatorTok{\}}

    \KeywordTok{private} \DataTypeTok{void} \FunctionTok{loadDashboardData}\OperatorTok{()} \OperatorTok{\{}
        \CommentTok{// Carga totales mensuales y calcula tendencia vs mes anterior}
\NormalTok{        executorService}\OperatorTok{.}\FunctionTok{execute}\OperatorTok{(()} \OperatorTok{{-}\textgreater{}} \OperatorTok{\{}
            \BuiltInTok{List}\OperatorTok{\textless{}}\NormalTok{MonthlyTotal}\OperatorTok{\textgreater{}}\NormalTok{ rawTotals }\OperatorTok{=}\NormalTok{ repository}\OperatorTok{.}\FunctionTok{getMonthlyTotalsSync}\OperatorTok{(}\DecValTok{3}\OperatorTok{);}
            \CommentTok{// Procesa y postea a monthlyExpenses y trendResult}
        \OperatorTok{\});}
    \OperatorTok{\}}

    \KeywordTok{public} \BuiltInTok{String} \FunctionTok{getRecomendacionDashboard}\OperatorTok{()} \OperatorTok{\{}
        \ControlFlowTok{return}\NormalTok{ FinancialTrendHelper}\OperatorTok{.}\FunctionTok{generateRecommendation}\OperatorTok{(}
\NormalTok{            estado}\OperatorTok{.}\FunctionTok{getEstado}\OperatorTok{(),}\NormalTok{ presupuesto}\OperatorTok{,}\NormalTok{ total}\OperatorTok{);}
    \OperatorTok{\}}

    \CommentTok{// CRUD: insertar, actualizar, borrar}
\OperatorTok{\}}
\end{Highlighting}
\end{Shaded}

\textbf{WishlistViewModel.java} - Wishlist con evaluación financiera
reactiva:

\begin{Shaded}
\begin{Highlighting}[]
\AttributeTok{@HiltViewModel}
\KeywordTok{public} \KeywordTok{class}\NormalTok{ WishlistViewModel }\KeywordTok{extends}\NormalTok{ AndroidViewModel }\OperatorTok{\{}
    \KeywordTok{private} \DataTypeTok{final}\NormalTok{ MediatorLiveData}\OperatorTok{\textless{}}\BuiltInTok{List}\OperatorTok{\textless{}}\NormalTok{WishlistItemUI}\OperatorTok{\textgreater{}\textgreater{}}\NormalTok{ wishlistEvaluada}\OperatorTok{;}

    \KeywordTok{public} \FunctionTok{WishlistViewModel}\OperatorTok{(}\NormalTok{Application app}\OperatorTok{,}\NormalTok{ WishlistRepository repo}\OperatorTok{,}
\NormalTok{            UserRepository userRepository}\OperatorTok{,}\NormalTok{ LanzamientoRepository repoLanzamiento}\OperatorTok{,}
\NormalTok{            GastoRepository gastoRepo}\OperatorTok{)} \OperatorTok{\{}

        \CommentTok{// Combina 3 fuentes de datos: wishlist, presupuesto y gasto actual}
\NormalTok{        wishlistEvaluada}\OperatorTok{.}\FunctionTok{addSource}\OperatorTok{(}\NormalTok{wishListSource}\OperatorTok{,}\NormalTok{ lista }\OperatorTok{{-}\textgreater{}}
            \FunctionTok{recalcularEvaluacion}\OperatorTok{(}\NormalTok{lista}\OperatorTok{,}\NormalTok{ presupuesto}\OperatorTok{,}\NormalTok{ gasto}\OperatorTok{));}
\NormalTok{        wishlistEvaluada}\OperatorTok{.}\FunctionTok{addSource}\OperatorTok{(}\NormalTok{presupuestoSource}\OperatorTok{,}\NormalTok{ presupuesto }\OperatorTok{{-}\textgreater{}}
            \FunctionTok{recalcularEvaluacion}\OperatorTok{(}\NormalTok{wishlist}\OperatorTok{,}\NormalTok{ presupuesto}\OperatorTok{,}\NormalTok{ gasto}\OperatorTok{));}
\NormalTok{        wishlistEvaluada}\OperatorTok{.}\FunctionTok{addSource}\OperatorTok{(}\NormalTok{gastoSource}\OperatorTok{,}\NormalTok{ gasto }\OperatorTok{{-}\textgreater{}}
            \FunctionTok{recalcularEvaluacion}\OperatorTok{(}\NormalTok{wishlist}\OperatorTok{,}\NormalTok{ presupuesto}\OperatorTok{,}\NormalTok{ gasto}\OperatorTok{));}
    \OperatorTok{\}}

    \KeywordTok{public} \DataTypeTok{void} \FunctionTok{comprarJuego}\OperatorTok{(}\NormalTok{WishlistEntity juego}\OperatorTok{,} \DataTypeTok{double}\NormalTok{ precioFinal}\OperatorTok{)} \OperatorTok{\{}
        \BuiltInTok{Date}\NormalTok{ fechaLanzamiento }\OperatorTok{=}\NormalTok{ FechaUtils}\OperatorTok{.}\FunctionTok{parseFecha}\OperatorTok{(}\NormalTok{juego}\OperatorTok{.}\FunctionTok{getFechaLanzamiento}\OperatorTok{());}
        \ControlFlowTok{if} \OperatorTok{(}\NormalTok{fechaLanzamiento }\OperatorTok{!=} \KeywordTok{null} \OperatorTok{\&\&}\NormalTok{ fechaLanzamiento}\OperatorTok{.}\FunctionTok{after}\OperatorTok{(}\KeywordTok{new} \BuiltInTok{Date}\OperatorTok{()))} \OperatorTok{\{}
            \CommentTok{// Es lanzamiento futuro {-}\textgreater{} guardar en Lanzamientos}
\NormalTok{            repoLanzamiento}\OperatorTok{.}\FunctionTok{insertar}\OperatorTok{(}\KeywordTok{new} \FunctionTok{LanzamientoEntity}\OperatorTok{(}\KeywordTok{...}\OperatorTok{));}
        \OperatorTok{\}} \ControlFlowTok{else} \OperatorTok{\{}
            \CommentTok{// Ya lanzado {-}\textgreater{} crear Gasto y eliminar de wishlist}
\NormalTok{            gastoRepo}\OperatorTok{.}\FunctionTok{insertarGasto}\OperatorTok{(}\KeywordTok{new} \FunctionTok{GastoEntity}\OperatorTok{(}\KeywordTok{...}\OperatorTok{));}
\NormalTok{            repo}\OperatorTok{.}\FunctionTok{borrar}\OperatorTok{(}\NormalTok{juego}\OperatorTok{);}
        \OperatorTok{\}}
    \OperatorTok{\}}
\OperatorTok{\}}
\end{Highlighting}
\end{Shaded}

\textbf{GamesViewModel.java} - Catálogo de juegos desde RAWG API:

\begin{Shaded}
\begin{Highlighting}[]
\AttributeTok{@HiltViewModel}
\KeywordTok{public} \KeywordTok{class}\NormalTok{ GamesViewModel }\KeywordTok{extends}\NormalTok{ ViewModel }\OperatorTok{\{}
    \KeywordTok{private} \DataTypeTok{final}\NormalTok{ GamesRepository repository}\OperatorTok{;}
    \KeywordTok{private} \DataTypeTok{final}\NormalTok{ MutableLiveData}\OperatorTok{\textless{}}\BuiltInTok{List}\OperatorTok{\textless{}}\NormalTok{GameDto}\OperatorTok{\textgreater{}\textgreater{}}\NormalTok{ juegos}\OperatorTok{;}

    \KeywordTok{public} \DataTypeTok{void} \FunctionTok{cargarJuegosRecientes}\OperatorTok{()} \OperatorTok{\{}
\NormalTok{        juegosRecientesLiveData }\OperatorTok{=}\NormalTok{ repository}\OperatorTok{.}\FunctionTok{getRecentlyReleasedGames}\OperatorTok{();}
\NormalTok{        juegosRecientesObserver }\OperatorTok{=}\NormalTok{ lista }\OperatorTok{{-}\textgreater{}} \OperatorTok{\{}
\NormalTok{            juegos}\OperatorTok{.}\FunctionTok{setValue}\OperatorTok{(}\FunctionTok{procesarLista}\OperatorTok{(}\NormalTok{lista}\OperatorTok{));}
        \OperatorTok{\};}
    \OperatorTok{\}}

    \KeywordTok{private} \BuiltInTok{List}\OperatorTok{\textless{}}\NormalTok{GameDto}\OperatorTok{\textgreater{}} \FunctionTok{procesarLista}\OperatorTok{(}\BuiltInTok{List}\OperatorTok{\textless{}}\NormalTok{GameDto}\OperatorTok{\textgreater{}}\NormalTok{ lista}\OperatorTok{)} \OperatorTok{\{}
        \ControlFlowTok{for} \OperatorTok{(}\NormalTok{GameDto juego }\OperatorTok{:}\NormalTok{ lista}\OperatorTok{)} \OperatorTok{\{}
            \CommentTok{// Calcula precio estimado según rating}
            \DataTypeTok{double}\NormalTok{ precio }\OperatorTok{=}\NormalTok{ sistema}\OperatorTok{.}\FunctionTok{estimarPrecio}\OperatorTok{(}\KeywordTok{new} \FunctionTok{GameInfo}\OperatorTok{(}\NormalTok{juego}\OperatorTok{.}\FunctionTok{getName}\OperatorTok{(),}\NormalTok{ juego}\OperatorTok{.}\FunctionTok{getRating}\OperatorTok{()));}
\NormalTok{            juego}\OperatorTok{.}\FunctionTok{setPrecioEstimado}\OperatorTok{(}\NormalTok{precio}\OperatorTok{);}
            \CommentTok{// Genera texto de plataformas}
\NormalTok{            juego}\OperatorTok{.}\FunctionTok{setPlataformasTexto}\OperatorTok{(}\FunctionTok{extractPlataformas}\OperatorTok{(}\NormalTok{juego}\OperatorTok{));}
        \OperatorTok{\}}
        \ControlFlowTok{return}\NormalTok{ lista}\OperatorTok{;}
    \OperatorTok{\}}

    \KeywordTok{public} \DataTypeTok{void} \FunctionTok{buscarJuegosPaginados}\OperatorTok{(}\BuiltInTok{String}\NormalTok{ query}\OperatorTok{,} \DataTypeTok{boolean}\NormalTok{ reset}\OperatorTok{)} \OperatorTok{\{}
        \CommentTok{// Búsqueda con paginación: reset=true reemplaza, reset=false añade}
    \OperatorTok{\}}
\OperatorTok{\}}
\end{Highlighting}
\end{Shaded}

\hypertarget{actividades-principales}{%
\subsection{Actividades principales}\label{actividades-principales}}

\textbf{MainActivity.java} - Dashboard principal:

\begin{Shaded}
\begin{Highlighting}[]
\AttributeTok{@AndroidEntryPoint}
\KeywordTok{public} \KeywordTok{class}\NormalTok{ MainActivity }\KeywordTok{extends}\NormalTok{ BaseNavActivity }\OperatorTok{\{}
    \AttributeTok{@Inject}\NormalTok{ UserRepository userRepository}\OperatorTok{;}
    \AttributeTok{@Inject}\NormalTok{ GastoRepository gastoRepo}\OperatorTok{;}

    \KeywordTok{protected} \DataTypeTok{void} \FunctionTok{onCreate}\OperatorTok{(}\NormalTok{Bundle savedInstanceState}\OperatorTok{)} \OperatorTok{\{}
        \FunctionTok{setContentView}\OperatorTok{(}\NormalTok{R}\OperatorTok{.}\FunctionTok{layout}\OperatorTok{.}\FunctionTok{activity\_main}\OperatorTok{);}
        \FunctionTok{inicializarVistas}\OperatorTok{();}
        \FunctionTok{configurarDrawer}\OperatorTok{();}
        \FunctionTok{inicializarViewModel}\OperatorTok{();}
        \FunctionTok{cargarPresupuestoUsuario}\OperatorTok{();}
        \FunctionTok{configurarObservers}\OperatorTok{();}  \CommentTok{// Observa estado financiero, gastos, tendencias}
        \FunctionTok{observarMoneda}\OperatorTok{();}       \CommentTok{// LiveData de moneda desde Firestore}
        \FunctionTok{configurarMenu}\OperatorTok{();}
        \FunctionTok{setupBottomNavigation}\OperatorTok{(}\NormalTok{R}\OperatorTok{.}\FunctionTok{id}\OperatorTok{.}\FunctionTok{nav\_home}\OperatorTok{);}
    \OperatorTok{\}}

    \KeywordTok{private} \DataTypeTok{void} \FunctionTok{configurarObservers}\OperatorTok{()} \OperatorTok{\{}
\NormalTok{        gastoViewModel}\OperatorTok{.}\FunctionTok{getEstadoUI}\OperatorTok{().}\FunctionTok{observe}\OperatorTok{(}\KeywordTok{this}\OperatorTok{,}\NormalTok{ ui }\OperatorTok{{-}\textgreater{}} \OperatorTok{\{}
\NormalTok{            tvRecomendacion}\OperatorTok{.}\FunctionTok{setText}\OperatorTok{(}\NormalTok{ui}\OperatorTok{.}\FunctionTok{getMensaje}\OperatorTok{());}
\NormalTok{            indicadorEstado}\OperatorTok{.}\FunctionTok{setBackgroundColor}\OperatorTok{(}\FunctionTok{getColor}\OperatorTok{(}\NormalTok{ui}\OperatorTok{.}\FunctionTok{getColorRes}\OperatorTok{()));}
        \OperatorTok{\});}

\NormalTok{        gastoViewModel}\OperatorTok{.}\FunctionTok{getListaGastos}\OperatorTok{().}\FunctionTok{observe}\OperatorTok{(}\KeywordTok{this}\OperatorTok{,}\NormalTok{ gastos }\OperatorTok{{-}\textgreater{}} \OperatorTok{\{}
            \CommentTok{// Actualiza PieChart y totales}
        \OperatorTok{\});}

\NormalTok{        gastoViewModel}\OperatorTok{.}\FunctionTok{getMonthlyExpenses}\OperatorTok{().}\FunctionTok{observe}\OperatorTok{(}\KeywordTok{this}\OperatorTok{,}\NormalTok{ expenses }\OperatorTok{{-}\textgreater{}} \OperatorTok{\{}
            \FunctionTok{actualizarLineChart}\OperatorTok{(}\NormalTok{expenses}\OperatorTok{);}
        \OperatorTok{\});}

\NormalTok{        gastoViewModel}\OperatorTok{.}\FunctionTok{getTrendResult}\OperatorTok{().}\FunctionTok{observe}\OperatorTok{(}\KeywordTok{this}\OperatorTok{,}\NormalTok{ trend }\OperatorTok{{-}\textgreater{}} \OperatorTok{\{}
            \FunctionTok{actualizarIndicadorTendencia}\OperatorTok{(}\NormalTok{trend}\OperatorTok{);}
        \OperatorTok{\});}
    \OperatorTok{\}}
\OperatorTok{\}}
\end{Highlighting}
\end{Shaded}

\textbf{ListaJuegosActivity.java} - Catálogo con búsqueda y paginación:

\begin{Shaded}
\begin{Highlighting}[]
\AttributeTok{@AndroidEntryPoint}
\KeywordTok{public} \KeywordTok{class}\NormalTok{ ListaJuegosActivity }\KeywordTok{extends}\NormalTok{ BaseNavActivity }\OperatorTok{\{}
    \KeywordTok{private} \DataTypeTok{static} \DataTypeTok{final} \DataTypeTok{long}\NormalTok{ DEBOUNCE\_DELAY }\OperatorTok{=} \DecValTok{500}\OperatorTok{;} \CommentTok{//ms}

    \KeywordTok{protected} \DataTypeTok{void} \FunctionTok{onCreate}\OperatorTok{(}\NormalTok{Bundle savedInstanceState}\OperatorTok{)} \OperatorTok{\{}
        \CommentTok{// Intenta usar datos precargados desde SplashActivity}
\NormalTok{        AgenteGamerApplication app }\OperatorTok{=} \OperatorTok{(}\NormalTok{AgenteGamerApplication}\OperatorTok{)} \FunctionTok{getApplication}\OperatorTok{();}
        \BuiltInTok{List}\OperatorTok{\textless{}}\NormalTok{GameDto}\OperatorTok{\textgreater{}}\NormalTok{ preloadedGames }\OperatorTok{=}\NormalTok{ app}\OperatorTok{.}\FunctionTok{getPreloadedGames}\OperatorTok{();}
        \ControlFlowTok{if} \OperatorTok{(}\NormalTok{preloadedGames }\OperatorTok{!=} \KeywordTok{null} \OperatorTok{\&\&} \OperatorTok{!}\NormalTok{preloadedGames}\OperatorTok{.}\FunctionTok{isEmpty}\OperatorTok{())} \OperatorTok{\{}
\NormalTok{            adapter}\OperatorTok{.}\FunctionTok{setLista}\OperatorTok{(}\NormalTok{preloadedGames}\OperatorTok{);}
        \OperatorTok{\}}

        \CommentTok{// Búsqueda con debounce}
\NormalTok{        searchView}\OperatorTok{.}\FunctionTok{addTextChangedListener}\OperatorTok{(}\KeywordTok{new} \FunctionTok{TextWatcher}\OperatorTok{()} \OperatorTok{\{}
            \KeywordTok{public} \DataTypeTok{void} \FunctionTok{onTextChanged}\OperatorTok{(}\BuiltInTok{CharSequence}\NormalTok{ s}\OperatorTok{,} \KeywordTok{...}\OperatorTok{)} \OperatorTok{\{}
\NormalTok{                handler}\OperatorTok{.}\FunctionTok{removeCallbacks}\OperatorTok{(}\NormalTok{searchRunnable}\OperatorTok{);}
\NormalTok{                searchRunnable }\OperatorTok{=} \OperatorTok{()} \OperatorTok{{-}\textgreater{}}\NormalTok{ viewModel}\OperatorTok{.}\FunctionTok{buscarJuegosPaginados}\OperatorTok{(}\NormalTok{texto}\OperatorTok{,} \KeywordTok{true}\OperatorTok{);}
\NormalTok{                handler}\OperatorTok{.}\FunctionTok{postDelayed}\OperatorTok{(}\NormalTok{searchRunnable}\OperatorTok{,}\NormalTok{ DEBOUNCE\_DELAY}\OperatorTok{);}
            \OperatorTok{\}}
        \OperatorTok{\});}

        \CommentTok{// Paginación infinita}
\NormalTok{        recyclerView}\OperatorTok{.}\FunctionTok{addOnScrollListener}\OperatorTok{(}\KeywordTok{new} \FunctionTok{OnScrollListener}\OperatorTok{()} \OperatorTok{\{}
            \KeywordTok{public} \DataTypeTok{void} \FunctionTok{onScrolled}\OperatorTok{(}\NormalTok{RecyclerView rv}\OperatorTok{,} \DataTypeTok{int}\NormalTok{ dx}\OperatorTok{,} \DataTypeTok{int}\NormalTok{ dy}\OperatorTok{)} \OperatorTok{\{}
                \ControlFlowTok{if} \OperatorTok{(}\NormalTok{lm}\OperatorTok{.}\FunctionTok{findLastVisibleItemPosition}\OperatorTok{()} \OperatorTok{\textgreater{}=}\NormalTok{ adapter}\OperatorTok{.}\FunctionTok{getItemCount}\OperatorTok{()} \OperatorTok{{-}} \DecValTok{3}\OperatorTok{)} \OperatorTok{\{}
\NormalTok{                    viewModel}\OperatorTok{.}\FunctionTok{buscarJuegosPaginados}\OperatorTok{(}\NormalTok{queryActual}\OperatorTok{,} \KeywordTok{false}\OperatorTok{);}
                \OperatorTok{\}}
            \OperatorTok{\}}
        \OperatorTok{\});}

        \CommentTok{// Click en juego {-}\textgreater{} añadir a wishlist}
\NormalTok{        adapter}\OperatorTok{.}\FunctionTok{setOnJuegoClickListener}\OperatorTok{((}\NormalTok{juego}\OperatorTok{,}\NormalTok{ precioEstimado}\OperatorTok{)} \OperatorTok{{-}\textgreater{}} \OperatorTok{\{}
\NormalTok{            wishlistViewModel}\OperatorTok{.}\FunctionTok{insertar}\OperatorTok{(}\KeywordTok{new} \FunctionTok{WishlistEntity}\OperatorTok{(}\KeywordTok{...}\OperatorTok{));}
        \OperatorTok{\});}
    \OperatorTok{\}}
\OperatorTok{\}}
\end{Highlighting}
\end{Shaded}

\textbf{BaseNavActivity.java} - Clase base para navegación inferior:

\begin{Shaded}
\begin{Highlighting}[]
\KeywordTok{public} \KeywordTok{abstract} \KeywordTok{class}\NormalTok{ BaseNavActivity }\KeywordTok{extends}\NormalTok{ AppCompatActivity }\OperatorTok{\{}
    \KeywordTok{private} \DataTypeTok{int}\NormalTok{ mSelectedItemId }\OperatorTok{=} \OperatorTok{{-}}\DecValTok{1}\OperatorTok{;}

    \KeywordTok{protected} \DataTypeTok{void} \FunctionTok{setupBottomNavigation}\OperatorTok{(}\DataTypeTok{int}\NormalTok{ selectedItemId}\OperatorTok{)} \OperatorTok{\{}
\NormalTok{        mSelectedItemId }\OperatorTok{=}\NormalTok{ selectedItemId}\OperatorTok{;}
\NormalTok{        BottomNavigationView bottomNav }\OperatorTok{=} \FunctionTok{findViewById}\OperatorTok{(}\NormalTok{R}\OperatorTok{.}\FunctionTok{id}\OperatorTok{.}\FunctionTok{bottomNavigationView}\OperatorTok{);}
\NormalTok{        bottomNav}\OperatorTok{.}\FunctionTok{setSelectedItemId}\OperatorTok{(}\NormalTok{selectedItemId}\OperatorTok{);}

\NormalTok{        bottomNav}\OperatorTok{.}\FunctionTok{setOnItemSelectedListener}\OperatorTok{(}\NormalTok{item }\OperatorTok{{-}\textgreater{}} \OperatorTok{\{}
            \ControlFlowTok{if} \OperatorTok{(}\NormalTok{item}\OperatorTok{.}\FunctionTok{getItemId}\OperatorTok{()} \OperatorTok{==}\NormalTok{ selectedItemId}\OperatorTok{)} \ControlFlowTok{return} \KeywordTok{true}\OperatorTok{;}
            \FunctionTok{navigateToActivity}\OperatorTok{(}\NormalTok{item}\OperatorTok{.}\FunctionTok{getItemId}\OperatorTok{());}
            \ControlFlowTok{return} \KeywordTok{true}\OperatorTok{;}
        \OperatorTok{\});}
    \OperatorTok{\}}

    \AttributeTok{@Override}
    \KeywordTok{protected} \DataTypeTok{void} \FunctionTok{onResume}\OperatorTok{()} \OperatorTok{\{}
        \KeywordTok{super}\OperatorTok{.}\FunctionTok{onResume}\OperatorTok{();}
        \CommentTok{// Resetea la selección al volver de otra actividad}
        \ControlFlowTok{if} \OperatorTok{(}\NormalTok{mSelectedItemId }\OperatorTok{!=} \OperatorTok{{-}}\DecValTok{1}\OperatorTok{)} \OperatorTok{\{}
\NormalTok{            BottomNavigationView bottomNav }\OperatorTok{=} \FunctionTok{findViewById}\OperatorTok{(}\NormalTok{R}\OperatorTok{.}\FunctionTok{id}\OperatorTok{.}\FunctionTok{bottomNavigationView}\OperatorTok{);}
            \ControlFlowTok{if} \OperatorTok{(}\NormalTok{bottomNav }\OperatorTok{!=} \KeywordTok{null} \OperatorTok{\&\&}\NormalTok{ bottomNav}\OperatorTok{.}\FunctionTok{getSelectedItemId}\OperatorTok{()} \OperatorTok{!=}\NormalTok{ mSelectedItemId}\OperatorTok{)} \OperatorTok{\{}
\NormalTok{                bottomNav}\OperatorTok{.}\FunctionTok{setSelectedItemId}\OperatorTok{(}\NormalTok{mSelectedItemId}\OperatorTok{);}
            \OperatorTok{\}}
        \OperatorTok{\}}
    \OperatorTok{\}}

    \KeywordTok{private} \DataTypeTok{void} \FunctionTok{navigateToActivity}\OperatorTok{(}\DataTypeTok{int}\NormalTok{ itemId}\OperatorTok{)} \OperatorTok{\{}
        \CommentTok{// Mapea itemId {-}\textgreater{} Activity y navega con CLEAR\_TOP | SINGLE\_TOP}
\NormalTok{        Intent intent }\OperatorTok{=} \KeywordTok{new} \FunctionTok{Intent}\OperatorTok{(}\KeywordTok{this}\OperatorTok{,}\NormalTok{ targetActivity}\OperatorTok{);}
\NormalTok{        intent}\OperatorTok{.}\FunctionTok{setFlags}\OperatorTok{(}\NormalTok{Intent}\OperatorTok{.}\FunctionTok{FLAG\_ACTIVITY\_CLEAR\_TOP} \OperatorTok{|}\NormalTok{ Intent}\OperatorTok{.}\FunctionTok{FLAG\_ACTIVITY\_SINGLE\_TOP}\OperatorTok{);}
        \FunctionTok{startActivity}\OperatorTok{(}\NormalTok{intent}\OperatorTok{);}
    \OperatorTok{\}}
\OperatorTok{\}}
\end{Highlighting}
\end{Shaded}

\hypertarget{utilidades}{%
\subsection{Utilidades}\label{utilidades}}

\textbf{MoneyUtils.java} - Formateo y conversión de moneda:

\begin{Shaded}
\begin{Highlighting}[]
\KeywordTok{public} \KeywordTok{class}\NormalTok{ MoneyUtils }\OperatorTok{\{}
    \KeywordTok{private} \DataTypeTok{static} \DataTypeTok{final} \BuiltInTok{Map}\OperatorTok{\textless{}}\BuiltInTok{String}\OperatorTok{,} \BuiltInTok{Double}\OperatorTok{\textgreater{}}\NormalTok{ EXCHANGE\_RATES }\OperatorTok{=} \KeywordTok{new} \BuiltInTok{HashMap}\OperatorTok{\textless{}\textgreater{}();}
    \DataTypeTok{static} \OperatorTok{\{}
\NormalTok{        EXCHANGE\_RATES}\OperatorTok{.}\FunctionTok{put}\OperatorTok{(}\StringTok{"EUR"}\OperatorTok{,} \FloatTok{1.0}\OperatorTok{);}
\NormalTok{        EXCHANGE\_RATES}\OperatorTok{.}\FunctionTok{put}\OperatorTok{(}\StringTok{"USD"}\OperatorTok{,} \FloatTok{1.08}\OperatorTok{);}
\NormalTok{        EXCHANGE\_RATES}\OperatorTok{.}\FunctionTok{put}\OperatorTok{(}\StringTok{"GBP"}\OperatorTok{,} \FloatTok{0.86}\OperatorTok{);}
    \OperatorTok{\}}

    \KeywordTok{public} \DataTypeTok{static} \BuiltInTok{String} \FunctionTok{format}\OperatorTok{(}\DataTypeTok{double}\NormalTok{ cantidad}\OperatorTok{,} \BuiltInTok{String}\NormalTok{ currencyCode}\OperatorTok{)} \OperatorTok{\{}
        \DataTypeTok{double}\NormalTok{ convertida }\OperatorTok{=} \FunctionTok{convertirDesdeEUR}\OperatorTok{(}\NormalTok{cantidad}\OperatorTok{,}\NormalTok{ currencyCode}\OperatorTok{);}
        \BuiltInTok{NumberFormat}\NormalTok{ formato }\OperatorTok{=} \BuiltInTok{NumberFormat}\OperatorTok{.}\FunctionTok{getCurrencyInstance}\OperatorTok{(}\NormalTok{LOCALE\_ES}\OperatorTok{);}
\NormalTok{        formato}\OperatorTok{.}\FunctionTok{setCurrency}\OperatorTok{(}\BuiltInTok{Currency}\OperatorTok{.}\FunctionTok{getInstance}\OperatorTok{(}\NormalTok{currencyCode}\OperatorTok{));}
        \ControlFlowTok{return}\NormalTok{ formato}\OperatorTok{.}\FunctionTok{format}\OperatorTok{(}\NormalTok{convertida}\OperatorTok{);}
    \OperatorTok{\}}

    \KeywordTok{public} \DataTypeTok{static} \DataTypeTok{double} \FunctionTok{convertirDesdeEUR}\OperatorTok{(}\DataTypeTok{double}\NormalTok{ montoEUR}\OperatorTok{,} \BuiltInTok{String}\NormalTok{ monedaDestino}\OperatorTok{)} \OperatorTok{\{}
        \BuiltInTok{Double}\NormalTok{ tasa }\OperatorTok{=}\NormalTok{ EXCHANGE\_RATES}\OperatorTok{.}\FunctionTok{get}\OperatorTok{(}\NormalTok{monedaDestino}\OperatorTok{.}\FunctionTok{toUpperCase}\OperatorTok{());}
        \ControlFlowTok{return}\NormalTok{ tasa }\OperatorTok{!=} \KeywordTok{null} \OperatorTok{?}\NormalTok{ montoEUR }\OperatorTok{*}\NormalTok{ tasa }\OperatorTok{:}\NormalTok{ montoEUR}\OperatorTok{;}
    \OperatorTok{\}}
\OperatorTok{\}}
\end{Highlighting}
\end{Shaded}

\hypertarget{inyecciuxf3n-de-dependencias}{%
\subsection{Inyección de
dependencias}\label{inyecciuxf3n-de-dependencias}}

\textbf{DatabaseModule.java}:

\begin{Shaded}
\begin{Highlighting}[]
\AttributeTok{@Module}
\AttributeTok{@InstallIn}\OperatorTok{(}\NormalTok{SingletonComponent}\OperatorTok{.}\FunctionTok{class}\OperatorTok{)}
\KeywordTok{public} \KeywordTok{class}\NormalTok{ DatabaseModule }\OperatorTok{\{}
    \AttributeTok{@Provides} \AttributeTok{@Singleton}
    \KeywordTok{public} \DataTypeTok{static}\NormalTok{ AppDatabase }\FunctionTok{provideDatabase}\OperatorTok{(}\AttributeTok{@ApplicationContext} \BuiltInTok{Context}\NormalTok{ context}\OperatorTok{)} \OperatorTok{\{}
        \ControlFlowTok{return}\NormalTok{ Room}\OperatorTok{.}\FunctionTok{databaseBuilder}\OperatorTok{(}\NormalTok{context}\OperatorTok{,}\NormalTok{ AppDatabase}\OperatorTok{.}\FunctionTok{class}\OperatorTok{,} \StringTok{"agente\_gamer\_db"}\OperatorTok{)}
            \OperatorTok{.}\FunctionTok{fallbackToDestructiveMigration}\OperatorTok{().}\FunctionTok{build}\OperatorTok{();}
    \OperatorTok{\}}

    \AttributeTok{@Provides} \KeywordTok{public} \DataTypeTok{static}\NormalTok{ GastoDao }\FunctionTok{provideGastoDao}\OperatorTok{(}\NormalTok{AppDatabase db}\OperatorTok{)} \OperatorTok{\{}
        \ControlFlowTok{return}\NormalTok{ db}\OperatorTok{.}\FunctionTok{gastoDao}\OperatorTok{();}
    \OperatorTok{\}}
    \AttributeTok{@Provides} \KeywordTok{public} \DataTypeTok{static}\NormalTok{ WishlistDao }\FunctionTok{provideWishlistDao}\OperatorTok{(}\NormalTok{AppDatabase db}\OperatorTok{)} \OperatorTok{\{}
        \ControlFlowTok{return}\NormalTok{ db}\OperatorTok{.}\FunctionTok{wishlistDao}\OperatorTok{();}
    \OperatorTok{\}}
    \AttributeTok{@Provides} \KeywordTok{public} \DataTypeTok{static}\NormalTok{ LanzamientoDao }\FunctionTok{provideLanzamientoDao}\OperatorTok{(}\NormalTok{AppDatabase db}\OperatorTok{)} \OperatorTok{\{}
        \ControlFlowTok{return}\NormalTok{ db}\OperatorTok{.}\FunctionTok{lanzamientoDao}\OperatorTok{();}
    \OperatorTok{\}}
    \AttributeTok{@Provides} \AttributeTok{@Singleton}
    \KeywordTok{public} \DataTypeTok{static} \BuiltInTok{ExecutorService} \FunctionTok{provideExecutorService}\OperatorTok{()} \OperatorTok{\{}
        \ControlFlowTok{return} \BuiltInTok{Executors}\OperatorTok{.}\FunctionTok{newSingleThreadExecutor}\OperatorTok{();}
    \OperatorTok{\}}
\OperatorTok{\}}
\end{Highlighting}
\end{Shaded}

\textbf{RepositoryModule.java}:

\begin{Shaded}
\begin{Highlighting}[]
\AttributeTok{@Module}
\AttributeTok{@InstallIn}\OperatorTok{(}\NormalTok{SingletonComponent}\OperatorTok{.}\FunctionTok{class}\OperatorTok{)}
\KeywordTok{public} \KeywordTok{abstract} \KeywordTok{class}\NormalTok{ RepositoryModule }\OperatorTok{\{}
    \AttributeTok{@Binds} \AttributeTok{@Singleton}
    \KeywordTok{public} \KeywordTok{abstract}\NormalTok{ GastoRepository }\FunctionTok{bindGastoRepository}\OperatorTok{(}\NormalTok{GastoRepositoryImpl impl}\OperatorTok{);}

    \AttributeTok{@Binds} \AttributeTok{@Singleton}
    \KeywordTok{public} \KeywordTok{abstract}\NormalTok{ WishlistRepository }\FunctionTok{bindWishlistRepository}\OperatorTok{(}\NormalTok{WishlistRepositoryImpl impl}\OperatorTok{);}

    \AttributeTok{@Binds} \AttributeTok{@Singleton}
    \KeywordTok{public} \KeywordTok{abstract}\NormalTok{ LanzamientoRepository }\FunctionTok{bindLanzamientoRepository}\OperatorTok{(}\NormalTok{LanzamientoRepositoryImpl impl}\OperatorTok{);}

    \AttributeTok{@Binds} \AttributeTok{@Singleton}
    \KeywordTok{public} \KeywordTok{abstract}\NormalTok{ UserRepository }\FunctionTok{bindUserRepository}\OperatorTok{(}\NormalTok{UserRepositoryImpl impl}\OperatorTok{);}
\OperatorTok{\}}
\end{Highlighting}
\end{Shaded}

\hypertarget{workers}{%
\subsection{Workers}\label{workers}}

\textbf{SistemaFinancieroWorker.java} - Validación periódica de compras:

\begin{Shaded}
\begin{Highlighting}[]
\AttributeTok{@HiltWorker}
\KeywordTok{public} \KeywordTok{class}\NormalTok{ SistemaFinancieroWorker }\KeywordTok{extends}\NormalTok{ Worker }\OperatorTok{\{}
    \KeywordTok{public} \BuiltInTok{Result} \FunctionTok{doWork}\OperatorTok{()} \OperatorTok{\{}
        \BuiltInTok{List}\OperatorTok{\textless{}}\NormalTok{WishlistEntity}\OperatorTok{\textgreater{}}\NormalTok{ wishlist }\OperatorTok{=}\NormalTok{ wishlistRepo}\OperatorTok{.}\FunctionTok{getWishlistSync}\OperatorTok{();}
        \CommentTok{// Obtiene presupuesto y gasto actual de forma asíncrona}
        \CommentTok{// Valida cada juego contra el presupuesto}
        \ControlFlowTok{for} \OperatorTok{(}\NormalTok{WishlistEntity juego }\OperatorTok{:}\NormalTok{ wishlist}\OperatorTok{)} \OperatorTok{\{}
            \BuiltInTok{String}\NormalTok{ evaluacion }\OperatorTok{=}\NormalTok{ validatePurchaseUseCase}\OperatorTok{.}\FunctionTok{validate}\OperatorTok{(}\NormalTok{juego}\OperatorTok{,}\NormalTok{ presupuesto}\OperatorTok{,}\NormalTok{ totalGastado}\OperatorTok{);}
            \ControlFlowTok{if} \OperatorTok{(!}\NormalTok{evaluacion}\OperatorTok{.}\FunctionTok{equals}\OperatorTok{(}\StringTok{"Compra recomendada"}\OperatorTok{))} \OperatorTok{\{}
                \FunctionTok{lanzarNotificacionNoComprar}\OperatorTok{(}\NormalTok{juego}\OperatorTok{);}
                \ControlFlowTok{break}\OperatorTok{;}
            \OperatorTok{\}}
        \OperatorTok{\}}
        \CommentTok{// Notifica sobre lanzamientos próximos (0{-}7 días)}
        \FunctionTok{avisarLanzamientosProximos}\OperatorTok{(}\NormalTok{wishlist}\OperatorTok{);}
        \ControlFlowTok{return} \BuiltInTok{Result}\OperatorTok{.}\FunctionTok{success}\OperatorTok{();}
    \OperatorTok{\}}
\OperatorTok{\}}
\end{Highlighting}
\end{Shaded}

\newpage

\hypertarget{anexo-ii---opiniuxf3n-personal-sobre-la-fct}{%
\section{11.- Anexo II - Opinión personal sobre la
FCT}\label{anexo-ii---opiniuxf3n-personal-sobre-la-fct}}

Mi experiencia durante las prácticas en empresa ha sido extremadamente
enriquecedora y me ha permitido aplicar todos los conocimientos
adquiridos durante el ciclo formativo de Desarrollo de Aplicaciones
Multiplataforma.

Durante las FCT he podido trabajar con tecnologías y herramientas que
solo habíamos visto de forma teórica en clase, como Firebase, Room,
Dagger Hilt y APIs REST. La diferencia más notable entre el entorno
académico y el profesional es la importancia de las buenas prácticas:
arquitectura limpia, separación de responsabilidades, testing y
documentación.

Uno de los mayores aprendizajes ha sido comprender la importancia de la
planificación. En clase, los proyectos suelen ser más lineales y con
requisitos claros. En el mundo real, los requisitos evolucionan, surgen
bugs inesperados y hay que saber priorizar. Aprender a gestionar el
tiempo y a tomar decisiones técnicas con criterio ha sido fundamental.

También he valorado mucho la experiencia de trabajar con Git en un
entorno colaborativo, creando ramas, pull requests y siguiendo un flujo
de trabajo estructurado. Esto me ha enseñado la importancia de los
commits atómicos, los mensajes descriptivos y la revisión de código.

El proyecto Agente Gamer ha sido el resultado directo de todo este
aprendizaje. Ha sido un reto técnico importante integrar múltiples
tecnologías (Firebase, Room, Retrofit, MPAndroidChart, Hilt) en una
arquitectura coherente, pero el resultado final me enorgullece y
demuestra que el ciclo ha cumplido su objetivo: prepararme para el mundo
profesional del desarrollo de software.

En resumen, las FCT han sido la pieza que faltaba para conectar la
teoría con la práctica, y el proyecto final es la prueba de que estoy
preparado para dar el salto al mundo laboral como desarrollador.

\end{document}
